\documentclass{article}
\usepackage{anyfontsize}
\usepackage[UTF8]{ctex} % 如果需要支持中文
\usepackage{fontspec} 

\usepackage[a4paper, left=0.1in, right=0.1in, top=0.05in, bottom=0.05in]{geometry} % 设置四边页边距为0.1英寸{geometry} % 设置页边距为0.5英寸
\usepackage{enumitem} % 用于自定义列表
\usepackage{amsmath}  % 数学公式支持

\setlength{\parindent}{0pt} % 不缩进
\usepackage{etoolbox} % 用于列表操作
%调整类代码
\usepackage{comment}
\usepackage{tocloft} % 用于定制目录样式
\usepackage{hyperref} %不需要目录盒子注释这
\usepackage{ifthen}
\usepackage{tikz}
\usepackage{titletoc}
\usepackage{graphicx}
\usepackage{amssymb}  % 提供数学符号，包括\therefore
%目录设定
%快捷键 CMD + / 注释
%  \renewcommand{\cftdot}{} % 隐藏目录中的页码
%  \cftpagenumbersoff{section}
%  \cftpagenumbersoff{subsection}
%  \makeatletter
%  \renewcommand{\@pnumwidth}{0em}  % 设置页码宽度为0
%  \renewcommand{\@tocrmarg}{0em}   % 设置右边距为0
%  \makeatother
 
%\setCJKmainfont[
 %   Path = C:/USERS/11252/APPDATA/LOCAL/MICROSOFT/WINDOWS/FONTS/,
 %   UprightFont = {SourceHanSerifCN-Light-5.otf},
 %   BoldFont = {SourceHanSerifCN-Medium-6.otf},
%    Scale=1
%]{SourceHanSerifCN}

%\setmainfont{SourceHanSerif}  % 设置英文字体

% 处理冲突的命令
\let\oldbackepsilon\backepsilon
\let\backepsilon\relax
\let\oldeth\eth
\let\eth\relax
\let\olddigamma\digamma
\let\digamma\relax

\usepackage{float}
\usepackage[a4paper, left=0.1in, right=0.1in, top=0.05in, bottom=0.05in]{geometry} % 设置四边页边距为0.1英寸{geometry} % 设置页边距为0.5英寸
\usepackage{enumitem} % 用于自定义列表
\usepackage{amsmath}  % 数学公式支持

\setlength{\parindent}{0pt} % 不缩进
\usepackage{etoolbox} % 用于列表操作
%调整类代码
\usepackage{comment}
\usepackage{tocloft} % 用于定制目录 样式
\usepackage{hyperref} %不需要目录盒子注释这
\usepackage{ifthen}
\usepackage{tikz}
\usepackage{titletoc}
\usepackage{graphicx}
\usepackage{amssymb}  % 提供数学符号，包括\therefore
%目录设定
%快捷键 CMD + / 注释
%  \renewcommand{\cftdot}{} % 隐藏目录中的页码
%  \cftpagenumbersoff{section}
%  \cftpagenumbersoff{subsection}
%  \makeatletter
%  \renewcommand{\@pnumwidth}{0em}  % 设置页码宽度为0
%  \renewcommand{\@tocrmarg}{0em}   % 设置右边距为0
%  \makeatother
 


\usepackage{environ}

\newif\ifshowcontent  % 定义一个条件判断变量
\showcontenttrue    % 设置为 false，隐藏所有 question 环境的内容

\newcounter{question} % 题目计数器
\setcounter{question}{0}
\NewEnviron{question}{%
    \ifshowcontent
        \par\stepcounter{question}\textbf{\thequestion.} \BODY\par % 如果显示内容，则输出
    \fi
}

\NewEnviron{choices}{%
    \ifshowcontent
        \begin{enumerate}[label=\Alph*., leftmargin=*]
            \BODY
        \end{enumerate}\par
    \fi
}

\NewEnviron{correctanswer}{%
    \ifshowcontent
        \textbf{答案:}\BODY\par
    \fi
}



\newenvironment{explanation}
{\textbf{解析:}\par}
{\par}



% 新增考察方面命令
\newcommand{\checklist}[1]{
    \textbf{考察:} #1 \par % 在题目下方显示考察方面
    \addcontentsline{toc}{subsection}{题号\thequestion 考察: #1} % 将考察内容添加到目录
    \vspace{2em} % 添加1em的垂直空间
}

\graphicspath{{images/}} %统一


\begin{document}
 %\fontsize{22}{31}\selectfont
 \tableofcontents % 添加目录
\newpage
%\pagestyle{empty}






\section{考点1:对手方衍生品风险}
\setcounter{question}{0}

\begin{question}
    Which of the following statements best describes the benefit of netting risk exposures? The benefits of netting are realized when:
\end{question}

\begin{choices}
    \item Marked-to-market (mtm) values have high structural correlations for two trades
    \item Marked-to-market (mtm) values have opposite signs for two trades
    \item Expected exposure (EE) values are minimal
    \item Expected future exposure (EFE) values have zero correlation
\end{choices}

\begin{correctanswer}
    B
\end{correctanswer}

\begin{explanation}
    \textbf{A选项错误。}
    高结构相关性并不意味着可以实现净额结算的好处。即使两笔交易高度相关，如果它们的MtM值符号相同，净额结算也无法减少风险敞口。

    \textbf{B选项正确。}
    当两笔交易的MtM值具有相反的符号时，净额结算可以显著减少风险敞口。例如，如果一笔交易的价值为+100，另一笔为-80，净额结算后的敞口仅为20，而不是180。

    \textbf{C选项错误。}
    预期敞口（EE）值最小化并不是净额结算的必要条件。即使EE值较大，只要MtM值符号相反，净额结算仍然可以带来好处。

    \textbf{D选项错误。}
    预期未来敞口（EFE）值的相关性为零并不意味着可以实现净额结算的好处。净额结算主要依赖于MtM值的符号，而不是EFE值的相关性。
\end{explanation}

\checklist{净额结算的好处在于当MtM值符号相反时可以减少风险敞口}


\begin{question}
    XYZ Hedge Fund wants to get exposure to a high-yield pool of commercial loans without actually investing in the loans. It wants a leverage ratio of 8.0. If the hedge fund is willing to invest \$40 million in this investment, which credit derivative is best for them and what is their expected return given that the reference asset earns LIBOR plus 300 basis points, the counterparty earns LIBOR plus 160 basis points, and the required collateral earns 4.0\%?
\end{question}

\begin{choices}
    \item Total return swap with a 14.80\% return
    \item Asset-backed credit-linked note with a 12.40\% return
    \item Total return swap with a 12.40\% return
    \item Asset-backed credit-linked note with a 14.80\% return
\end{choices}

\begin{correctanswer}
    D
\end{correctanswer}

\begin{explanation}
    \textbf{A选项错误。}
    总收益互换（Total Return Swap）在这种情况下不是最佳选择，因为资产支持信用关联票据（Asset-backed Credit-linked Note）可以提供更高的回报。

    \textbf{B选项错误。}
    资产支持信用关联票据的回报计算错误。正确的计算应该考虑杠杆效应和利差。

    \textbf{C选项错误。}
    总收益互换的回报计算错误，且不是最佳选择。

    \textbf{D选项正确。}
    对冲基金应该选择资产支持信用关联票据。计算过程如下：
    \begin{itemize}
        \item 杠杆率为8.0，投资额为4000万美元，因此商业贷款池的名义价值为3.2亿美元
        \item 对冲基金将从4000万美元的抵押品中获得4.0\%的收益，即160万美元
        \item 对冲基金还将从整个3.2亿美元中获得140个基点的利差（300-160），即448万美元
        \item 对冲基金的总回报率为14.80\% [（160万美元 + 448万美元）/ 4000万美元]
    \end{itemize}
\end{explanation}

\checklist{资产支持信用关联票据的回报计算需要考虑杠杆效应、抵押品收益和利差}


\begin{question}
    A portfolio manager enters into a total return swap as the total return receiver. Under which of the following situations, would the portfolio manager be required to make a net outlay to the counterparty?
\end{question}

\begin{choices}
    \item If the transaction was initiated as a hedge. Then no outlay was required
    \item If there were a capital gain on the reference asset
    \item If the market value of the reference asset decreased significantly
    \item If the spread between the reference asset and the benchmark asset changed
\end{choices}

\begin{correctanswer}
    C
\end{correctanswer}

\begin{explanation}
    \textbf{A选项错误。}
    交易是否作为对冲工具发起与是否需要支付净额无关。总收益互换的支付义务取决于参考资产的表现，而不是交易的目的。

    \textbf{B选项错误。}
    如果参考资产出现资本利得，总收益接收方（投资组合经理）将收到支付，而不是需要支付。总收益接收方在资产升值时获得收益，在资产贬值时承担损失。

    \textbf{C选项正确。}
    当参考资产的市场价值显著下降时，总收益接收方需要向对手方支付净差额。这是因为总收益接收方同时面临信用风险和市场风险，必须承担参考资产价值下降带来的损失。

    \textbf{D选项错误。}
    参考资产与基准资产之间的利差变化本身并不会导致总收益接收方需要支付净额。支付义务主要取决于参考资产的总收益表现。
\end{explanation}

\checklist{总收益接收方在参考资产价值下降时需要支付净差额}


\begin{question}
    Which of the following statements is not correct regarding total return swaps (TRS)?
\end{question}

\begin{choices}
    \item A TRS is designed to mirror the return on an underlying asset like a loan, stock, or even a portfolio of assets
    \item The payer pays any depreciation in the underlying asset to the receiver
    \item The payer pays any dividends or interest received to the receiver
    \item The receiver is creating a synthetic long position in the underlying
\end{choices}

\begin{correctanswer}
    B
\end{correctanswer}

\begin{explanation}
    \textbf{A选项错误。}
    这个说法是正确的。总收益互换（TRS）的设计目的就是复制标的资产（如贷款、股票或资产组合）的收益。

    \textbf{B选项正确。}
    这个说法是错误的。在总收益互换中，支付方（payer）只负责支付标的资产的升值部分以及相关的股息或利息。而接收方（receiver）需要向支付方支付标的资产的贬值部分。总收益互换同时转移了信用风险和市场风险。

    \textbf{C选项错误。}
    这个说法是正确的。支付方确实需要将收到的任何股息或利息支付给接收方，这是总收益互换的基本机制之一。

    \textbf{D选项错误。}
    这个说法是正确的。接收方通过总收益互换实际上创建了一个对标的资产的合成多头头寸，因为接收方获得了标的资产的所有收益（包括升值和收入），同时也承担了所有损失。
\end{explanation}

\checklist{总收益互换中支付方不负责支付标的资产的贬值部分}


\begin{question}
    In a total return swap, the:
\end{question}

\begin{choices}
    \item Total return payer bears the risk of holding the reference asset
    \item Total return payer transfers the risk of holding the reference asset to the total return receiver
    \item Total return receiver transfers the risk of holding the reference asset to the total return payer
    \item Credit risk of the buyer and seller cancel each other out
\end{choices}

\begin{correctanswer}
    B
\end{correctanswer}

\begin{explanation}
    \textbf{A选项错误。}
    总收益支付方（payer）并不承担持有参考资产的风险。相反，他们通过总收益互换将风险转移给了总收益接收方。

    \textbf{B选项正确。}
    在总收益互换中，总收益支付方将参考资产的资本利得和损失转移给了总收益接收方。总收益互换是一种信用衍生产品，它将标的资产（如贷款、债券或资产组合）的总收益与LIBOR加上信用价差进行交换。总收益包括利息收入和资产价值变动。总收益互换不仅可以管理信用风险，还可以管理市场风险。

    \textbf{C选项错误。}
    总收益接收方（receiver）并不将风险转移给支付方，而是从支付方那里接收风险。接收方通过支付LIBOR加上价差来获得参考资产的总收益。

    \textbf{D选项错误。}
    买卖双方的信用风险并不会相互抵消。总收益互换中的信用风险主要存在于交易对手方之间，而不是相互抵消。
\end{explanation}

\checklist{总收益互换中支付方将参考资产的风险转移给接收方}

\begin{question}
    Risk Averse Bank (RAB) has made a loan of USD 120 million at 7.5\% per annum. RAB wants to enter into a total return swap under which it will pay the interest on the loan plus the change in the mark-to-market value of the loan and in exchange, RAB will get LIBOR + 40bps. Settlement payments are made annually. What is the cash flow for RAB on the first settlement data if the mark-to-market value of the loan falls by 2.5\% and LIBOR is 5.5\%?
\end{question}

\begin{choices}
    \item Net inflow of USD 0.4 million
    \item Net outflow of USD 0.4 million
    \item Net inflow of USD 2.0 million
    \item Net outflow of USD 2.0 million
\end{choices}

\begin{correctanswer}
    A
\end{correctanswer}

\begin{explanation}
    \textbf{A选项正确。}
    计算过程如下：
    \begin{itemize}
        \item RAB在第一年年底获得贷款利息：120百万 × 7.5\% = 9百万美元
        \item 贷款价值下降：120百万 × 2.5\% = 3百万美元
        \item RAB需要支付的总收益：9百万 - 3百万 = 6百万美元
        \item RAB收到的LIBOR + 40bps：120百万 × (5.5\% + 0.4\%) = 7.08百万美元
        \item 净流入：7.08百万 - 6百万 = 0.4百万美元
    \end{itemize}

    \textbf{B选项错误。}
    这个答案计算错误。RAB收到的LIBOR + 40bps（7.08百万美元）大于需要支付的总收益（6百万美元），因此是净流入而不是净流出。

    \textbf{C选项错误。}
    这个答案计算错误。净流入应该是0.4百万美元，而不是2.0百万美元。

    \textbf{D选项错误。}
    这个答案计算错误。RAB收到的LIBOR + 40bps（7.08百万美元）大于需要支付的总收益（6百万美元），因此是净流入而不是净流出。
\end{explanation}

\checklist{总收益互换中净现金流等于收到的LIBOR加价差减去支付的总收益}


\begin{question}
    The credit exposure of an interest rate swap differs from that of a bond in that:

    I. The swap can be terminated by novation

    II. The principal amount of the swap is not at risk

    III. Swaps benefit from higher recovery rates

    IV. The full coupon amounts in the swap are not at risk
\end{question}

\begin{choices}
    \item I and III
    \item II and IV
    \item II, III and IV
    \item I, II, III and IV
\end{choices}

\begin{correctanswer}
    B
\end{correctanswer}

\begin{explanation}
    \textbf{A选项错误。}
    这个选项包含了错误的说法。互换的回收率并不一定高于债券，而且通过更新终止（novation）并不是互换信用风险区别于债券的主要特征。

    \textbf{B选项正确。}
    利率互换的信用风险与债券有两个主要区别：
    \begin{itemize}
        \item 互换的本金金额不面临风险，因为名义本金不需要交换
        \item 互换中只有固定利率和浮动利率之间的净差额需要交换，而不是全部的息票金额
    \end{itemize}
    债券是资产负债表上的资产，其信用风险就是贷款金额。而利率互换的信用风险只是相同本金债券的一小部分，因为：
    \begin{itemize}
        \item 互换只有在合约价值为正时才会导致损失（而债券总是有正价值）
        \item 互换市场中的交易对手通常有净额结算协议，可以进一步降低信用风险
    \end{itemize}

    \textbf{C选项错误。}
    这个选项包含了错误的说法。互换的回收率并不一定高于债券。

    \textbf{D选项错误。}
    这个选项包含了错误的说法。通过更新终止（novation）和更高的回收率并不是互换信用风险区别于债券的主要特征。
\end{explanation}

\checklist{利率互换的信用风险小于债券，因为本金不面临风险且只有净差额需要交换}


\begin{question}
    An investor can gain the exposure and return of an underlying loan by:
\end{question}

\begin{choices}
    \item Being the credit risk buyer in a total return swap
    \item Being the credit risk seller in a total return swap
    \item Entering into an interest-rate swap with a BB credit
    \item Pledging compensating balances with the lender bank
\end{choices}

\begin{correctanswer}
    A
\end{correctanswer}

\begin{explanation}
    \textbf{A选项正确。}
    在总收益互换中，信用风险卖方（总收益支付方）将标的资产的总收益支付给信用风险买方（总收益接收方）。通过成为总收益互换中的信用风险买方，投资者可以获得标的贷款的风险敞口和收益，而不需要实际持有该贷款。

    \textbf{B选项错误。}
    作为总收益互换中的信用风险卖方，投资者实际上是转移了标的贷款的风险和收益，而不是获得这些风险和收益。信用风险卖方需要支付标的资产的总收益，并收取LIBOR加上价差。

    \textbf{C选项错误。}
    与BB级信用的对手方进行利率互换并不能获得标的贷款的风险敞口和收益。利率互换只涉及利率的交换，不涉及信用风险的转移。

    \textbf{D选项错误。}
    向贷款银行提供补偿性余额并不能获得标的贷款的风险敞口和收益。补偿性余额只是作为贷款的担保，并不涉及风险和收益的转移。
\end{explanation}

\checklist{总收益互换中的信用风险买方可以获得标的贷款的风险敞口和收益}


\begin{question}
    A standard synthetic CDO (basket credit default swap) references a portfolio of twelve (12) individual corporate names. Assume the following:
    • The total reference notional (basket notional) is X, and the term is Y years
    • The reference notional per individual reference credit name is X/12 (i.e. equal weight per name)
    • The default correlations between the individual reference credit names are all equal to one (1.0)
    • The single-name credit default swap (CDS) spread for each individual reference credit name is 120 basis points, for a term of Y years
    • The assumed recovery rate on default for all individual reference credits is zero in all cases
    • The synthetic CDO comprises two tranches, a 40% junior tranche priced at a spread J, and a 60% senior tranche priced at spread S
    All else held constant, if the default correlations between the individual reference credit names are reduced from 1.0 to 0.6, what is the effect on the relationship between the junior tranche spread J and the senior tranche spread S?
\end{question}

\begin{choices}
    \item The relationship remains the same
    \item S increases relative to J
    \item J increases relative to S
    \item The effect cannot be determined given the data supplied
\end{choices}

\begin{correctanswer}
    C
\end{correctanswer}

\begin{explanation}
    \textbf{A选项错误。}
    当违约相关性从1.0降低到0.6时，层级之间的关系会发生变化。相关性变化会影响不同层级的风险分配。

    \textbf{B选项错误。}
    当违约相关性降低时，高级层（senior tranche）的利差不会相对于次级层（junior tranche）增加。相反，次级层的利差会增加。

    \textbf{C选项正确。}
    当违约相关性从1.0降低到0.6时：
    \begin{itemize}
        \item 在相关性为1.0时，只有两种可能的结果：要么所有发行人同时违约，要么没有发行人违约
        \item 当相关性降低到0.6时，会出现部分发行人违约的情况
        \item 这些部分违约的情况对次级层的影响大于对高级层的影响
        \item 虽然整体平均违约率保持不变，但次级层承担了更多的风险
        \item 因此，次级层的利差相对于高级层会增加，因为整个组合的总利差必须保持在120个基点
    \end{itemize}

    \textbf{D选项错误。}
    根据提供的数据，可以确定违约相关性降低对层级利差关系的影响。次级层的利差会相对于高级层增加。
\end{explanation}

\checklist{违约相关性降低时，次级层利差相对于高级层增加}


\begin{question}
    Blackstone Credit, Inc., made a loan to a small start-up firm. The firm grew rapidly, and it appeared that Blackstone had made a good credit decision. However, the firm grew too fast and could not sustain the growth. It eventually failed. Blackstone had initially estimated its exposure at default to be \$1,500,000. Because of the firm's rapid growth and resulting increases in the line of credit, Blackstone ultimately lost \$1,850,000. In terms of credit risk, this is an example of:
\end{question}

\begin{choices}
    \item Default on payment for goods or services already rendered
    \item A more severe loss than expected due to a ratings downgrade by a rating agency
    \item A more severe loss than expected due to a greater than expected exposure at the time of a default
    \item A more severe loss than expected due to a lower than expected recovery at the time of a default
\end{choices}

\begin{correctanswer}
    C
\end{correctanswer}

\begin{explanation}
    \textbf{A选项错误。}
    这个案例涉及的是贷款违约，而不是已提供商品或服务的付款违约。Blackstone提供的是信贷资金，而不是商品或服务。

    \textbf{B选项错误。}
    这个案例中没有提到评级机构降级的情况。作为一家小型初创企业，该公司可能根本没有信用评级。

    \textbf{C选项正确。}
    这个案例展示了由于违约时的风险敞口大于最初估计而导致的损失超出预期的情况：
    \begin{itemize}
        \item Blackstone最初估计的违约风险敞口为150万美元
        \item 由于公司快速增长，信贷额度增加
        \item 最终损失达到185万美元，比预期多出35万美元
        \item 这种损失是由于违约时的实际风险敞口大于最初估计造成的
    \end{itemize}

    \textbf{D选项错误。}
    这个案例中没有提到回收率的问题。损失超出预期是由于风险敞口增加，而不是回收率低于预期。
\end{explanation}

\checklist{违约时的实际风险敞口大于预期导致损失超出预期}


\begin{question}
    ABC Company, domiciled in a country with a strong legal system, is applying for a USD 120,000 loan with an annual interest rate of 6% to be used exclusively on expanding its business and not to repay current debt. The loan will be secured by the company's factory buildings, which have an appraised value of USD 250,000. The company has an annual operating profit of USD 4,000. Which of the following statements relating to the credit risk of this company can most likely be inferred from the given data?
\end{question}

\begin{choices}
    \item The company's capacity to pay is low
    \item The company's willingness to pay is low
    \item The loan exhibits a high loss given default
    \item The loan's exposure at default will increase
\end{choices}

\begin{correctanswer}
    A
\end{correctanswer}

\begin{explanation}
    \textbf{A选项正确。}
    从给定数据可以推断出公司的偿付能力很低：
    \begin{itemize}
        \item 贷款金额为12万美元，年利率为6%
        \item 年利息支出为7,200美元（12万 × 6%）
        \item 公司年营业利润仅为4,000美元
        \item 营业利润不足以支付利息支出，更不用说偿还本金
        \item 这表明公司的偿付能力（capacity to pay）很低
    \end{itemize}

    \textbf{B选项错误。}
    从给定数据无法推断出公司的还款意愿（willingness to pay）很低。公司申请贷款用于业务扩张而不是偿还现有债务，这表明公司有良好的还款意愿。

    \textbf{C选项错误。}
    从给定数据无法推断出贷款在违约时会有高损失。抵押物（工厂建筑）的评估价值为25万美元，远高于贷款金额，这表明违约时的损失可能较低。

    \textbf{D选项错误。}
    从给定数据无法推断出违约时的风险敞口会增加。贷款金额是固定的12万美元，没有信息表明这个金额会随时间增加。
\end{explanation}

\checklist{公司的偿付能力很低，因为营业利润不足以支付利息支出}


\begin{question}
    An endowment fund has sold default protection on the most senior tranche of a CDO. If the default correlation between assets held in the CDO decreases sharply from the correlation used in pricing the CDO tranches, assuming everything else is unchanged, how will the position of the endowment fund be impacted?
\end{question}

\begin{choices}
    \item It will either increase or decrease, depending on the pricing model used and the market conditions
    \item It will gain significant value, since the probability of exercising the protection falls
    \item It will lose significant value, since the protection will gain value
    \item It will neither gain nor lose value, since only expected default losses matter and correlation does not affect expected default losses
\end{choices}

\begin{correctanswer}
    B
\end{correctanswer}

\begin{explanation}
    \textbf{A选项错误。}
    违约相关性降低对最高层级（senior tranche）的影响是确定的，不会因为定价模型或市场条件而改变方向。相关性降低总是会增加最高层级的价值。

    \textbf{B选项正确。}
    当违约相关性降低时，最高层级的价值会增加，原因如下：
    \begin{itemize}
        \item 高相关性意味着如果一个资产违约，CDO中的大量其他资产也会违约
        \item 低相关性意味着如果一个资产违约，对其他资产的违约概率影响很小
        \item 随着相关性降低，发生足够多违约以超过最高层级信用增强的概率也会降低
        \item 因此，出售最高层级保护的投资者（捐赠基金）将获得收益
    \end{itemize}

    \textbf{C选项错误。}
    这个说法与实际情况相反。违约相关性降低会增加最高层级的价值，而不是降低其价值。

    \textbf{D选项错误。}
    违约相关性确实会影响CDO的价值，特别是对最高层级的影响。相关性变化会影响违约的集中度，从而影响不同层级的风险。
\end{explanation}

\checklist{违约相关性降低会增加CDO最高层级的价值}


\begin{question}
    When considering counterparty credit risk, which of the following financial products has the largest outstanding notional amount in the marketplace?
\end{question}

\begin{choices}
    \item Credit default swaps
    \item Foreign exchange forwards
    \item Interest rate swaps
    \item Repos and reverse repos
\end{choices}

\begin{correctanswer}
    C
\end{correctanswer}

\begin{explanation}
    \textbf{A选项错误。}
    信用违约互换（CDS）虽然是一个重要的场外衍生品市场，但其名义金额远小于利率互换市场。CDS主要用于信用风险转移，市场规模相对较小。

    \textbf{B选项错误。}
    外汇远期合约是重要的场外衍生品，但其名义金额也远小于利率互换市场。外汇市场虽然交易活跃，但单个合约的规模通常较小。

    \textbf{C选项正确。}
    利率互换是场外衍生品市场中规模最大的产品：
    \begin{itemize}
        \item 利率互换占据了场外衍生品市场的大部分份额
        \item 利率互换被广泛用于利率风险管理和资产负债管理
        \item 由于利率风险管理的普遍需求，利率互换的名义金额远大于其他金融产品
    \end{itemize}

    \textbf{D选项错误。}
    回购和逆回购虽然也是重要的金融产品，但它们属于证券融资交易，而不是场外衍生品。它们的市场规模小于利率互换市场。
\end{explanation}

\checklist{利率互换是场外衍生品市场中规模最大的产品}


\begin{question}
    All of the following could help to reduce the credit exposure on a set of derivative transactions except:
\end{question}

\begin{choices}
    \item Daily mark-to-market of transactions
    \item Netting agreements with the counterparty
    \item Collateral and other credit enhancements
    \item Early termination agreements
\end{choices}

\begin{correctanswer}
    D
\end{correctanswer}

\begin{explanation}
    \textbf{A选项错误。}
    每日盯市（Daily mark-to-market）是减少信用风险敞口的有效方法：
    \begin{itemize}
        \item 通过每日重新评估交易价值，可以及时发现风险敞口的变化
        \item 盯市机制允许及时调整保证金要求
        \item 有助于提前发现潜在的信用问题
    \end{itemize}

    \textbf{B选项错误。}
    净额结算协议（Netting agreements）是减少信用风险敞口的重要工具：
    \begin{itemize}
        \item 允许将多个交易的正负敞口相互抵消
        \item 显著降低总体信用风险敞口
        \item 是场外衍生品市场常用的风险管理工具
    \end{itemize}

    \textbf{C选项错误。}
    抵押品和其他信用增强措施是减少信用风险敞口的有效方法：
    \begin{itemize}
        \item 提供额外的信用保护
        \item 降低违约时的损失
        \item 包括现金抵押、证券抵押等多种形式
    \end{itemize}

    \textbf{D选项正确。}
    提前终止协议（Early termination agreements）是减少信用风险敞口效果最差的方法：
    \begin{itemize}
        \item 这些协议通常在交易对手信用质量下降时触发
        \item 此时可能已经难以有效执行终止条款
        \item 其他交易对手可能已经采取类似行动，导致市场流动性下降
    \end{itemize}
\end{explanation}

\checklist{提前终止协议是减少信用风险敞口效果最差的方法}


\begin{question}
    Which of the following statements regarding the International Swaps and Derivatives Association (ISDA) Master Agreement is false?
\end{question}

\begin{choices}
    \item The ISDA Master Agreement is a contract between two parties that governs derivative transactions and credit controls
    \item An ISDA Master Agreement is prepared for every derivatives transaction that occurs between two parties
    \item The ISDA Master Agreement forms a single contract under which multiple derivative contracts can be managed
    \item The ISDA Master Agreement specifies termination events that each pair of counterparties must accept
\end{choices}

\begin{correctanswer}
    B
\end{correctanswer}

\begin{explanation}
    \textbf{A选项错误。}
    这个说法是正确的。ISDA主协议确实是双方之间的合同，用于管理衍生品交易和信用控制：
    \begin{itemize}
        \item 规定了交易双方的权利和义务
        \item 包含了信用风险控制条款
        \item 是场外衍生品交易的基础法律文件
    \end{itemize}

    \textbf{B选项正确。}
    这个说法是错误的。ISDA主协议的特点：
    \begin{itemize}
        \item 不是为每笔交易单独准备一份协议
        \item 是一个框架协议，可以管理多个衍生品合同
        \item 一旦签署，后续交易只需通过确认书（confirmation）确认具体条款
    \end{itemize}

    \textbf{C选项错误。}
    这个说法是正确的。ISDA主协议确实是一个单一合同，可以管理多个衍生品合同：
    \begin{itemize}
        \item 提供了标准化的法律框架
        \item 允许在同一个协议下进行多笔交易
        \item 简化了交易管理流程
    \end{itemize}

    \textbf{D选项错误。}
    这个说法是正确的。ISDA主协议确实规定了交易双方必须接受的终止事件：
    \begin{itemize}
        \item 包括违约事件
        \item 包括终止事件
        \item 规定了终止后的处理程序
    \end{itemize}
\end{explanation}

\checklist{ISDA主协议是一个框架协议，可以管理多个衍生品合同，而不是为每笔交易单独准备}


\begin{question}
    Which of the following features is least likely a benefit of collateralization?
\end{question}

\begin{choices}
    \item Reduces capital requirements
    \item Allows for more competitive pricing of counterparty risk
    \item Reduces market, operational, and liquidity risk
    \item Reduces credit exposure
\end{choices}

\begin{correctanswer}
    C
\end{correctanswer}

\begin{explanation}
    \textbf{A选项错误。}
    抵押品确实可以降低资本要求：
    \begin{itemize}
        \item 通过提供额外的信用保护
        \item 降低风险加权资产
        \item 提高资本效率
    \end{itemize}

    \textbf{B选项错误。}
    抵押品确实允许对交易对手风险进行更具竞争力的定价：
    \begin{itemize}
        \item 降低信用风险溢价
        \item 使交易条件更具吸引力
        \item 提高市场竞争力
    \end{itemize}

    \textbf{C选项正确。}
    抵押品实际上会增加而不是降低市场、操作和流动性风险：
    \begin{itemize}
        \item 市场风险：当抵押品价值下降时，可能导致敞口部分或完全无抵押
        \item 操作风险：在获取抵押品时可能遇到法律障碍
        \item 流动性风险：在需要时可能难以以公平市场价格出售抵押品
    \end{itemize}

    \textbf{D选项错误。}
    抵押品确实可以降低信用风险敞口：
    \begin{itemize}
        \item 提供额外的信用保护
        \item 降低违约时的损失
        \item 改善信用风险状况
    \end{itemize}
\end{explanation}

\checklist{抵押品会增加市场、操作和流动性风险，而不是降低这些风险}


\begin{question}
    An investor can gain the exposure and return of an underlying loan by:
\end{question}

\begin{choices}
    \item Being the credit risk buyer in a total return swap
    \item Being the credit risk seller in a total return swap
    \item Entering into an interest-rate swap with a BB credit
    \item Pledging compensating balances with the lender bank
\end{choices}

\begin{correctanswer}
    A
\end{correctanswer}

\begin{explanation}
    \textbf{A选项正确。}
    在总收益互换中，信用风险买方（总收益接收方）可以获得标的贷款的风险敞口和收益：
    \begin{itemize}
        \item 信用风险卖方（总收益支付方）将标的资产的总收益支付给买方
        \item 买方获得标的贷款的所有收益，包括利息收入和资本利得/损失
        \item 买方承担标的贷款的信用风险和市场风险
    \end{itemize}

    \textbf{B选项错误。}
    作为总收益互换中的信用风险卖方，投资者实际上是转移了标的贷款的风险和收益，而不是获得这些风险和收益：
    \begin{itemize}
        \item 卖方需要支付标的资产的总收益
        \item 卖方收取LIBOR加上价差
        \item 卖方转移了信用风险和市场风险
    \end{itemize}

    \textbf{C选项错误。}
    与BB级信用的对手方进行利率互换并不能获得标的贷款的风险敞口和收益：
    \begin{itemize}
        \item 利率互换只涉及利率的交换
        \item 不涉及信用风险的转移
    \end{itemize}

    \textbf{D选项错误。}
    向贷款银行提供补偿性余额并不能获得标的贷款的风险敞口和收益：
    \begin{itemize}
        \item 补偿性余额只是作为贷款的担保
        \item 不涉及风险和收益的转移
    \end{itemize}
\end{explanation}

\checklist{总收益互换中的信用风险买方可以获得标的贷款的风险敞口和收益}


\begin{question}
    Which of the following statements is not correct regarding total return swaps (TRS)?
\end{question}

\begin{choices}
    \item A TRS is designed to mirror the return on an underlying asset like a loan, stock, or even a portfolio of assets
    \item The payer pays any depreciation in the underlying asset to the receiver
    \item The payer pays any dividends or interest received to the receiver
    \item The receiver is creating a synthetic long position in the underlying
\end{choices}

\begin{correctanswer}
    B
\end{correctanswer}

\begin{explanation}
    \textbf{A选项错误。}
    这个说法是正确的。总收益互换（TRS）的设计目的就是复制标的资产（如贷款、股票或资产组合）的收益：
    \begin{itemize}
        \item 包括利息或股息收入
        \item 包括资本利得或损失
        \item 完全反映标的资产的表现
    \end{itemize}

    \textbf{B选项正确。}
    这个说法是错误的。在总收益互换中：
    \begin{itemize}
        \item 支付方（payer）只负责支付标的资产的升值部分以及相关的股息或利息
        \item 接收方（receiver）需要向支付方支付标的资产的贬值部分
        \item 总收益互换同时转移了信用风险和市场风险
    \end{itemize}

    \textbf{C选项错误。}
    这个说法是正确的。支付方确实需要将收到的任何股息或利息支付给接收方，这是总收益互换的基本机制之一。

    \textbf{D选项错误。}
    这个说法是正确的。接收方通过总收益互换实际上创建了一个对标的资产的合成多头头寸：
    \begin{itemize}
        \item 获得标的资产的所有收益
        \item 承担标的资产的所有损失
        \item 相当于持有标的资产
    \end{itemize}
\end{explanation}

\checklist{总收益互换中支付方不负责支付标的资产的贬值部分}


\begin{question}
    In pricing a first-to-default credit basket swap, which of the following is true, all else being equal?
\end{question}

\begin{choices}
    \item The lower the correlation between the assets of the basket, the lower the premium
    \item The lower the correlation between the assets of the basket, the higher the premium
    \item The higher the correlation between the assets of the basket, the higher the premium
    \item The correlation between the assets has no impact in the premium of a first-to-default credit basket swap
\end{choices}

\begin{correctanswer}
    B
\end{correctanswer}

\begin{explanation}
    \textbf{A选项错误。}
    这个说法与实际情况相反。在first-to-default信用篮子互换中，资产相关性越低，保费越高。

    \textbf{B选项正确。}
    在first-to-default信用篮子互换中，资产相关性越低，保费越高，原因如下：
    \begin{itemize}
        \item first-to-default互换在篮子中任何实体首次违约时触发
        \item 当资产相关性较低时，篮子面临更大的违约风险
        \item 例如，如果篮子包含来自不同行业的资产，则篮子面临每个行业的违约风险
        \item 如果篮子只包含来自同一行业的资产，则相关性较高，违约风险较低
    \end{itemize}

    \textbf{C选项错误。}
    这个说法与实际情况相反。在first-to-default信用篮子互换中，资产相关性越高，保费越低。

    \textbf{D选项错误。}
    资产相关性确实会影响first-to-default信用篮子互换的保费：
    \begin{itemize}
        \item 当资产相关性为负时，first-to-default CDS保费更高
        \item 当资产相关性为正且为1时，first-to-default CDS与second-to-default CDS的保费相同
    \end{itemize}
\end{explanation}

\checklist{first-to-default信用篮子互换中，资产相关性越低，保费越高}



\begin{question}
    In a total return swap, the:
\end{question}

\begin{choices}
    \item Total return payer bears the risk of holding the reference asset
    \item Total return payer transfers the risk of holding the reference asset to the total return receiver
    \item Total return receiver transfers the risk of holding the reference asset to the total return payer
    \item Credit risk of the buyer and seller cancel each other out
\end{choices}

\begin{correctanswer}
    B
\end{correctanswer}

\begin{explanation}
    \textbf{A选项错误。}
    总收益支付方（payer）并不承担持有参考资产的风险。相反，他们通过总收益互换将风险转移给了总收益接收方。

    \textbf{B选项正确。}
    在总收益互换中，总收益支付方将参考资产的风险转移给总收益接收方：
    \begin{itemize}
        \item 支付方将参考资产的资本利得和损失转移给接收方
        \item 接收方获得参考资产的所有收益，包括利息收入和资本利得/损失
        \item 接收方承担参考资产的信用风险和市场风险
    \end{itemize}

    \textbf{C选项错误。}
    总收益接收方（receiver）并不将风险转移给支付方，而是从支付方那里接收风险。接收方通过支付LIBOR加上价差来获得参考资产的总收益。

    \textbf{D选项错误。}
    买卖双方的信用风险并不会相互抵消。总收益互换中的信用风险主要存在于交易对手方之间，而不是相互抵消。
\end{explanation}

\checklist{总收益互换中支付方将参考资产的风险转移给接收方}



\begin{question}
    ABC Company, domiciled in a country with a strong legal system, is applying for a USD 120,000 loan with an annual interest rate of 6\% to be used exclusively on expanding its business and not to repay current debt. The loan will be secured by the company's factory buildings, which have an appraised value of USD 250,000. The company has an annual operating profit of USD 4,000. Which of the following statements relating to the credit risk of this company can most likely be inferred from the given data?
\end{question}

\begin{choices}
    \item The company's capacity to pay is low
    \item The company's willingness to pay is low
    \item The loan exhibits a high loss given default
    \item The loan's exposure at default will increase
\end{choices}

\begin{correctanswer}
    A
\end{correctanswer}

\begin{explanation}
    \textbf{A选项正确。}
    从给定数据可以推断出公司的偿付能力很低：
    \begin{itemize}
        \item 贷款金额为12万美元，年利率为6\%
        \item 年利息支出为7,200美元（12万 × 6\%）
        \item 公司年营业利润仅为4,000美元
        \item 营业利润不足以支付利息支出，更不用说偿还本金
        \item 这表明公司的偿付能力（capacity to pay）很低
    \end{itemize}

    \textbf{B选项错误。}
    从给定数据无法推断出公司的还款意愿（willingness to pay）很低：
    \begin{itemize}
        \item 公司申请贷款用于业务扩张而不是偿还现有债务
        \item 这表明公司有良好的还款意愿
        \item 没有信息表明公司有逃避债务的意图
    \end{itemize}

    \textbf{C选项错误。}
    从给定数据无法推断出贷款在违约时会有高损失：
    \begin{itemize}
        \item 抵押物（工厂建筑）的评估价值为25万美元
        \item 抵押物价值远高于贷款金额
        \item 这表明违约时的损失可能较低
    \end{itemize}

    \textbf{D选项错误。}
    从给定数据无法推断出违约时的风险敞口会增加：
    \begin{itemize}
        \item 贷款金额是固定的12万美元
        \item 没有信息表明这个金额会随时间增加
    \end{itemize}
\end{explanation}

\checklist{公司的偿付能力很低，因为营业利润不足以支付利息支出}


\begin{question}
    Lin Ping is valuing a 1-year credit default swap (CDS) contract which will pay the buyer 80\% of the face value of a bond issued by Xiao Corp, immediately after a default by Xiao. To purchase this CDS, the buyer will pay the CDS spread, which is a percentage of the face value, once at the end of the year. Lin estimates that the risk-neutral default probability for Xiao is 6\% per year. The risk-free rate is 4\% per year. Assuming defaults can only occur halfway through the year and that the accrued premium is paid immediately after a default, what is the estimate for the CDS spread?
\end{question}

\begin{choices}
    \item 480 basis points
    \item 485 basis points
    \item 490 basis points
    \item 500 basis points
\end{choices}

\begin{correctanswer}
    B
\end{correctanswer}

\begin{explanation}
    \textbf{A选项错误。}
    这个答案计算错误。CDS利差的计算需要考虑违约概率、回收率和无风险利率。

    \textbf{B选项正确。}
    CDS利差的计算过程如下：
    \begin{itemize}
        \item 违约概率（PD）= 6\%
        \item 回收率（RR）= 20\%（1 - 80\%）
        \item 无风险利率（r）= 4\%
        \item 违约时间点 = 0.5年
        \item 预期损失 = PD × (1 - RR) = 6\% × 80\% = 4.8\%
        \item 考虑时间价值：4.8\% / (1 + 4\% × 0.5) = 4.85\%
        \item 因此，CDS利差为485个基点
    \end{itemize}

    \textbf{C选项错误。}
    这个答案计算错误。CDS利差的计算需要考虑违约概率、回收率和无风险利率。

    \textbf{D选项错误。}
    这个答案计算错误。CDS利差的计算需要考虑违约概率、回收率和无风险利率。
\end{explanation}

\checklist{CDS利差等于预期损失除以考虑时间价值的因子}



\newpage


\section{考点2:对手方风险延申}
\setcounter{question}{0}


\begin{question}
    Which of the following two transactions increases counterparty credit exposure?

    I. Selling a forward contract to the counterparty

    II. Selling a call option to the counterparty
\end{question}

\begin{choices}
    \item I only
    \item II only
    \item Both
    \item Neither
\end{choices}

\begin{correctanswer}
    A
\end{correctanswer}

\begin{explanation}
    \textbf{A选项正确。}
    只有出售远期合约会增加交易对手信用风险敞口：
    \begin{itemize}
        \item 远期合约的出售会产生交易对手信用风险
        \item 风险敞口取决于交易对手在到期日的履约能力
        \item 如果交易对手违约，卖方可能面临损失
    \end{itemize}

    \textbf{B选项错误。}
    出售看涨期权不会增加交易对手信用风险敞口：
    \begin{itemize}
        \item 期权卖方在交易开始时已经收取了期权费
        \item 期权卖方的义务是确定的，不依赖于交易对手的履约
        \item 交易对手违约不会对卖方造成负面影响
    \end{itemize}

    \textbf{C选项错误。}
    这个选项错误地认为两种交易都会增加信用风险敞口。实际上，只有远期合约的出售会增加信用风险敞口。

    \textbf{D选项错误。}
    这个选项错误地认为两种交易都不会增加信用风险敞口。实际上，远期合约的出售会增加信用风险敞口。
\end{explanation}

\checklist{只有出售远期合约会增加交易对手信用风险敞口}

\begin{question}
    Which of the following frictions represents an adverse selection problem?
\end{question}

\begin{choices}
    \item Investor and mortgagor
    \item Originator and arranger
    \item Servicer and rating agency
    \item Servicer and mortgagor
\end{choices}

\begin{correctanswer}
    B
\end{correctanswer}

\begin{explanation}
    \textbf{A选项错误。}
    投资者和抵押贷款人之间不存在典型的逆向选择问题：
    \begin{itemize}
        \item 投资者通过证券化产品间接持有抵押贷款
        \item 投资者主要面临的是信息不对称问题
        \item 不是典型的逆向选择情况
    \end{itemize}

    \textbf{B选项正确。}
    发起人（Originator）和承销人（Arranger）之间存在逆向选择问题：
    \begin{itemize}
        \item 发起人拥有关于借款人质量的更多信息
        \item 发起人可能保留高质量的抵押贷款
        \item 承销人可能收到低质量的抵押贷款（"柠檬"）
        \item 这是典型的逆向选择问题
    \end{itemize}

    \textbf{C选项错误。}
    服务商和评级机构之间不存在典型的逆向选择问题：
    \begin{itemize}
        \item 服务商负责贷款的管理和收款
        \item 评级机构负责评估证券化产品的信用风险
        \item 两者之间是服务关系，不是逆向选择
    \end{itemize}

    \textbf{D选项错误。}
    服务商和抵押贷款人之间不存在典型的逆向选择问题：
    \begin{itemize}
        \item 服务商负责贷款的管理和收款
        \item 抵押贷款人是借款人
        \item 两者之间是服务关系，不是逆向选择
    \end{itemize}
\end{explanation}

\checklist{发起人和承销人之间存在逆向选择问题}


\begin{question}
    Which of the following methods of mitigating counterparty risk is most likely to generate systemic risk?
\end{question}

\begin{choices}
    \item Netting
    \item Collateral
    \item Hedging
    \item Central counterparties
\end{choices}

\begin{correctanswer}
    D
\end{correctanswer}

\begin{explanation}
    \textbf{A选项错误。}
    净额结算（Netting）不会产生系统性风险：
    \begin{itemize}
        \item 净额结算只是将多个交易的正负敞口相互抵消
        \item 风险分散在多个交易对手之间
        \item 不会导致风险的集中
    \end{itemize}

    \textbf{B选项错误。}
    抵押品（Collateral）不会产生系统性风险：
    \begin{itemize}
        \item 抵押品是双边交易中的信用保护措施
        \item 风险分散在多个交易对手之间
        \item 不会导致风险的集中
    \end{itemize}

    \textbf{C选项错误。}
    对冲（Hedging）不会产生系统性风险：
    \begin{itemize}
        \item 对冲是将风险转移给其他市场参与者
        \item 风险分散在多个交易对手之间
        \item 不会导致风险的集中
    \end{itemize}

    \textbf{D选项正确。}
    中央对手方（Central counterparties）最可能产生系统性风险：
    \begin{itemize}
        \item 中央对手方将交易对手风险集中在少数几个机构
        \item 如果其中一个中央对手方失败，可能对整个金融系统造成重大冲击
        \item 风险的集中可能导致系统性风险
    \end{itemize}
\end{explanation}

\checklist{中央对手方将风险集中在少数几个机构，最可能产生系统性风险}


\begin{question}
    Liz Parker is a junior quantitative analyst who is preparing a report dealing with credit migration. An excerpt of her report contains the following statements:
    I. Future default probability will likely increase over time, especially for periods far into the future.
    II. When computing the default probability of a counterparty under a risk-neutral measure, we need to first determine the actual default probability.
    Which of Parker's statements is (are) correct?
\end{question}

\begin{choices}
    \item I only
    \item II only
    \item Both I and II
    \item Neither I nor II
\end{choices}

\begin{correctanswer}
    D
\end{correctanswer}

\begin{explanation}
    \textbf{A选项错误。}
    陈述I是错误的。未来的违约概率实际上会随着时间推移而降低：
    \begin{itemize}
        \item 违约更可能发生在较早的时间点
        \item 随着时间推移，违约已经发生的可能性增加
        \item 因此，未来的违约概率会降低
    \end{itemize}

    \textbf{B选项错误。}
    陈述II是错误的。在风险中性测度下计算违约概率时：
    \begin{itemize}
        \item 需要计算市场隐含的理论概率
        \item 实际违约概率适用于实际（历史）测度
        \item 不需要先确定实际违约概率
    \end{itemize}

    \textbf{C选项错误。}
    两个陈述都是错误的，因此这个选项不正确。

    \textbf{D选项正确。}
    两个陈述都是错误的：
    \begin{itemize}
        \item 陈述I错误：未来违约概率会随时间降低
        \item 陈述II错误：风险中性测度下使用市场隐含概率
    \end{itemize}
\end{explanation}

\checklist{两个陈述都是错误的：未来违约概率会随时间降低，风险中性测度下使用市场隐含概率}


\begin{question}
    Which of the following statements regarding counterparty credit risk is most accurate?
\end{question}

\begin{choices}
    \item Counterparty risk is unilateral
    \item Over-the-counter (OTC) derivatives contain less counterparty risk than exchange-traded derivatives because the counterparty is known
    \item The precise future value of the contract is uncertain, but the counterparties are aware of whether the future value will be positive or negative
    \item Counterparty risk is typically associated with counterparty default prior to the settlement rather than default during the settlement process
\end{choices}

\begin{correctanswer}
    D
\end{correctanswer}

\begin{explanation}
    \textbf{A选项错误。}
    交易对手风险不是单边的，而是双边的：
    \begin{itemize}
        \item 双方都不知道合同的最终价值
        \item 双方都不知道是否会盈利或亏损
        \item 双方都面临交易对手违约的风险
    \end{itemize}

    \textbf{B选项错误。}
    场外衍生品的交易对手风险实际上大于交易所交易的衍生品：
    \begin{itemize}
        \item 交易所交易的衍生品的交易对手是交易所
        \item 交易所通过保证金制度有效降低了交易对手风险
        \item 场外衍生品的交易对手风险更大
    \end{itemize}

    \textbf{C选项错误。}
    交易对手不仅不知道合同的确切未来价值，也不知道未来价值是正是负：
    \begin{itemize}
        \item 合同价值取决于未来的市场条件
        \item 双方都无法预知未来价值的方向
        \item 这是交易对手风险的本质特征
    \end{itemize}

    \textbf{D选项正确。}
    交易对手风险主要与结算前的违约有关：
    \begin{itemize}
        \item 虽然违约可能发生在结算前和结算期间
        \item 但交易对手风险主要关注结算前的违约风险
        \item 这是交易对手风险的主要表现形式
    \end{itemize}
\end{explanation}

\checklist{交易对手风险主要与结算前的违约有关}


\begin{question}
    Which of the following statements regarding lending and counterparty risk is true?
\end{question}

\begin{choices}
    \item Counterparty risk is defined as the possibility that an obligator will default on an outstanding loan
    \item Lending risk is defined as the possibility that either party to a derivative transaction will fail to meet its obligations
    \item Counterparty risk is very different from and more complex than lending risk
    \item Counterparty risk and lending risk are the same
\end{choices}

\begin{correctanswer}
    C
\end{correctanswer}

\begin{explanation}
    \textbf{A选项错误。}
    这个定义实际上是贷款风险的定义，而不是交易对手风险的定义：
    \begin{itemize}
        \item 贷款风险是指借款人可能违约的可能性
        \item 交易对手风险的定义与此不同
        \item 这个选项混淆了两种风险的定义
    \end{itemize}

    \textbf{B选项错误。}
    这个定义实际上是交易对手风险的定义，而不是贷款风险的定义：
    \begin{itemize}
        \item 交易对手风险是指衍生品交易中任何一方可能违约的可能性
        \item 贷款风险的定义与此不同
        \item 这个选项混淆了两种风险的定义
    \end{itemize}

    \textbf{C选项正确。}
    交易对手风险与贷款风险有很大区别，且更复杂：
    \begin{itemize}
        \item 交易对手风险是双向的，双方都可能违约
        \item 交易对手风险涉及未来不确定的现金流
        \item 交易对手风险需要考虑市场风险和信用风险的相互作用
        \item 交易对手风险的管理更复杂
    \end{itemize}

    \textbf{D选项错误。}
    交易对手风险和贷款风险是不同的：
    \begin{itemize}
        \item 贷款风险是单向的，只关注借款人的违约
        \item 交易对手风险是双向的，双方都可能违约
        \item 两种风险的特征和管理方法都不同
    \end{itemize}
\end{explanation}

\checklist{交易对手风险与贷款风险有很大区别，且更复杂}


\begin{question}
    Which of the following statements is not a motivation for pricing counterparty risk?
\end{question}

\begin{choices}
    \item Accurate pricing should only account for the cost of the trade
    \item Counterparty risk pricing should account for risk mitigants
    \item Best practices organize pricing responsibilities in the organization
    \item Pricing bilateral derivatives contracts
\end{choices}

\begin{correctanswer}
    A
\end{correctanswer}

\begin{explanation}
    \textbf{A选项正确。}
    这个说法不是交易对手风险定价的动机：
    \begin{itemize}
        \item 准确的定价不仅要考虑交易成本
        \item 还要考虑交易对手风险成本
        \item 只考虑交易成本是不完整的
    \end{itemize}

    \textbf{B选项错误。}
    考虑风险缓释措施是交易对手风险定价的重要动机：
    \begin{itemize}
        \item 风险缓释措施可以降低交易对手风险
        \item 需要在定价中反映风险缓释的效果
        \item 这是风险定价的重要组成部分
    \end{itemize}

    \textbf{C选项错误。}
    在组织中明确定价责任是交易对手风险定价的动机：
    \begin{itemize}
        \item 确保风险定价的准确性和一致性
        \item 明确各部门的职责
        \item 这是良好的风险管理实践
    \end{itemize}

    \textbf{D选项错误。}
    为双边衍生品合约定价是交易对手风险定价的动机：
    \begin{itemize}
        \item 双边衍生品合约涉及交易对手风险
        \item 需要在定价中反映这种风险
        \item 这是风险定价的重要应用
    \end{itemize}
\end{explanation}

\checklist{准确的定价不仅要考虑交易成本，还要考虑交易对手风险成本}


\begin{question}
    A risk manager at a large regional bank asks a group of junior analysts to study and report on the economic costs of OTC derivatives traded by the bank in relation to the bank's counterparty risk and funding risk. The manager instructs the analysts to use a hypothetical threshold that is higher for the counterparty than for the bank in assessing the bank's funding and collateral positions. Which of the following observations, from the perspective of the regional bank, would the analysts be correct to make?
\end{question}

\begin{choices}
    \item The bank has a funding cost but faces no counterparty risk when a derivative portfolio is in-the-money
    \item The bank faces no counterparty risk and no funding cost when a derivative portfolio is out-of-the-money
    \item The bank has a potential funding benefit when a derivative portfolio is in-the-money but not when it is out-of-the-money
    \item The bank faces counterparty risk whether a derivative portfolio is in-the-money or out-of-the-money
\end{choices}

\begin{correctanswer}
    B
\end{correctanswer}

\begin{explanation}
    \textbf{A选项错误。}
    当衍生品组合处于实值状态时：
    \begin{itemize}
        \item 银行确实面临融资成本
        \item 但银行仍然面临交易对手风险
        \item 因为交易对手可能违约
    \end{itemize}

    \textbf{B选项正确。}
    当衍生品组合处于虚值状态时：
    \begin{itemize}
        \item 银行不面临交易对手风险，因为交易对手违约不会造成损失
        \item 银行不面临融资成本，因为不需要提供抵押品
        \item 这是正确的观察
    \end{itemize}

    \textbf{C选项错误。}
    当衍生品组合处于实值状态时：
    \begin{itemize}
        \item 银行面临融资成本，而不是融资收益
        \item 因为需要提供抵押品
        \item 这个观察是错误的
    \end{itemize}

    \textbf{D选项错误。}
    银行面临的交易对手风险取决于衍生品组合的价值状态：
    \begin{itemize}
        \item 当组合处于虚值状态时，不面临交易对手风险
        \item 当组合处于实值状态时，才面临交易对手风险
        \item 这个观察是错误的
    \end{itemize}
\end{explanation}

\checklist{当衍生品组合处于虚值状态时，银行不面临交易对手风险和融资成本}






\newpage



\section{考点3:净额结算及相关方面}
\setcounter{question}{0}

\begin{question}
    Company ABC has great deal of transactions with counterparty XYZ. ABC has some contracts with XYZ covered under netting agreement A, has some contracts with XYZ covered under netting agreement B, and some contracts with XYZ with no netting agreement. With the following data, what is ABC's current counterparty credit exposure to XYZ?
\end{question}

\begin{choices}
    \item USD 8,500
    \item USD 7,200
    \item USD 5,800
    \item USD 10,200
\end{choices}

\begin{correctanswer}
    D
\end{correctanswer}

\begin{explanation}
    \textbf{A选项错误。}
    这个答案计算错误。需要分别计算三个部分的敞口：
    \begin{itemize}
        \item 净额协议A下的敞口
        \item 净额协议B下的敞口
        \item 无净额协议下的敞口
    \end{itemize}

    \textbf{B选项错误。}
    这个答案计算错误。需要分别计算三个部分的敞口：
    \begin{itemize}
        \item 净额协议A下的敞口
        \item 净额协议B下的敞口
        \item 无净额协议下的敞口
    \end{itemize}

    \textbf{C选项错误。}
    这个答案计算错误。需要分别计算三个部分的敞口：
    \begin{itemize}
        \item 净额协议A下的敞口
        \item 净额协议B下的敞口
        \item 无净额协议下的敞口
    \end{itemize}

    \textbf{D选项正确。}
    计算过程如下：
    \begin{itemize}
        \item 净额协议A下的敞口 = max(2,500 - 3,800 - 2,200, 0) = 0
        \item 净额协议B下的敞口 = max(6,500 - 800 - 2,800 + 4,500, 0) = 7,400
        \item 无净额协议下的敞口 = 2,800
        \item 总敞口 = 0 + 7,400 + 2,800 = 10,200
    \end{itemize}
\end{explanation}

\checklist{交易对手信用风险敞口等于各净额协议下的敞口加上无净额协议下的敞口}


\begin{question}
    Bank One has made a \$250 million loan to a software company at a fixed rate of 10 percent. The bank wants to hedge its exposure by entering into a Total Return Swap with a counterparty, Interloan Co., in which Bank One promises to pay the interest on the loan plus the change in the market value of the loan in exchange for LIBOR plus 50 basis points. If after one year, the market value of the loan has decreased by 4 percent and LIBOR is 9 percent, what will be the net obligation of Bank One?
\end{question}

\begin{choices}
    \item Net receipt of \$6.0 million
    \item Net payment of \$6.0 million
    \item Net receipt of \$6.5 million
    \item Net payment of \$6.5 million
\end{choices}

\begin{correctanswer}
    A
\end{correctanswer}

\begin{explanation}
    \textbf{A选项正确。}
    计算过程如下：
    \begin{itemize}
        \item 银行应收金额 = \$250 million × (9% + 0.5%) = \$23.75 million
        \item 银行应付金额 = \$250 million × 10% + \$250 million × (-4%) = \$15 million
        \item 净收入 = \$23.75 million - \$15 million = \$6.0 million
    \end{itemize}

    \textbf{B选项错误。}
    这个答案计算错误。银行是净收入方，而不是净支付方。

    \textbf{C选项错误。}
    这个答案计算错误。净收入金额应该是\$6.0 million，而不是\$6.5 million。

    \textbf{D选项错误。}
    这个答案计算错误。银行是净收入方，而不是净支付方，且金额计算错误。
\end{explanation}

\checklist{总收益互换中，银行的净义务等于应收金额减去应付金额}

\begin{question}
    A diversified portfolio of OTC derivatives with a single counterparty currently has a net mark-to-market of USD 25,000,000 and a gross mark-to-market (the sum of the value of all positive value positions minus the value of all negative value positions) of USD 90,000,000. Assuming there are no netting agreements in place with the counterparty, determine the current credit exposure to the counterparty.
\end{question}

\begin{choices}
    \item Less than or equal to USD 15,000,000
    \item Greater than USD 15,000,000 but less than or equal to USD 45,000,000
    \item Greater than USD 45,000,000 but less than USD 65,000,000
    \item Greater than USD 65,000,000
\end{choices}

\begin{correctanswer}
    C
\end{correctanswer}

\begin{explanation}
    \textbf{A选项错误。}
    这个答案计算错误。需要根据净值和总值计算正价值头寸的金额。

    \textbf{B选项错误。}
    这个答案计算错误。需要根据净值和总值计算正价值头寸的金额。

    \textbf{C选项正确。}
    计算过程如下：
    \begin{itemize}
        \item 设正价值头寸为A，负价值头寸为B
        \item 净值 = A - B = \$25 million
        \item 总值 = A + B = \$90 million
        \item 解得：A = \$57.5 million
        \item 因此，当前信用风险敞口为\$57.5 million
    \end{itemize}

    \textbf{D选项错误。}
    这个答案计算错误。需要根据净值和总值计算正价值头寸的金额。
\end{explanation}

\checklist{无净额协议时，信用风险敞口等于正价值头寸的金额}


\begin{question}
    Riggs Resources, LLC, (Riggs) is a commodity trading firm. Riggs has numerous trades outstanding with several counterparties; however, it is concerned with presettlement risk. In order to reduce presettlement risk (the risk that Riggs's counterparties would default before settlement), it would be most beneficial for Riggs to
\end{question}

\begin{choices}
    \item Have payment netting
    \item Have close-out netting
    \item Analyze potential losses as the sum of exposures
    \item Have netting but not set-off
\end{choices}

\begin{correctanswer}
    B
\end{correctanswer}

\begin{explanation}
    \textbf{A选项错误。}
    支付净额结算（payment netting）不能有效降低结算前风险：
    \begin{itemize}
        \item 支付净额结算主要降低结算风险和操作风险
        \item 不涉及交易对手违约的处理
        \item 不能有效降低结算前风险
    \end{itemize}

    \textbf{B选项正确。}
    终止净额结算（close-out netting）最能有效降低结算前风险：
    \begin{itemize}
        \item 在交易对手违约时，可以终止所有合同并进行净额结算
        \item 将多个交易的正负敞口相互抵消
        \item 显著降低结算前风险
    \end{itemize}

    \textbf{C选项错误。}
    将敞口简单相加分析潜在损失不能降低结算前风险：
    \begin{itemize}
        \item 这种方法忽略了净额结算的效果
        \item 会高估实际风险敞口
        \item 不能有效降低风险
    \end{itemize}

    \textbf{D选项错误。}
    净额结算和抵销（set-off）实际上是同义词：
    \begin{itemize}
        \item 两者都指将多个交易的正负敞口相互抵消
        \item 不能同时有净额结算而没有抵销
        \item 这个选项本身是矛盾的
    \end{itemize}
\end{explanation}

\checklist{终止净额结算最能有效降低结算前风险}


\begin{question}
    What does netting permit a firm to do?
\end{question}

\begin{choices}
    \item Pay only the losses due to unfavorable market prices
    \item Net all of its gains against all of its losses
    \item Net its transaction with a given counterparty against each other
    \item Pay the losses only if the counterpart defaults
\end{choices}

\begin{correctanswer}
    C
\end{correctanswer}

\begin{explanation}
    \textbf{A选项错误。}
    净额结算不是只支付不利市场价格造成的损失：
    \begin{itemize}
        \item 净额结算涉及所有交易的正负敞口
        \item 不仅限于市场风险造成的损失
        \item 这个描述不准确
    \end{itemize}

    \textbf{B选项错误。}
    净额结算不是将所有收益和损失相互抵消：
    \begin{itemize}
        \item 净额结算只适用于与同一交易对手的交易
        \item 不能跨交易对手进行净额结算
        \item 这个描述不准确
    \end{itemize}

    \textbf{C选项正确。}
    净额结算允许公司将与特定交易对手的交易相互抵消：
    \begin{itemize}
        \item 将应收交易对手的款项与应付交易对手的款项相抵消
        \item 减少实际支付金额
        \item 降低信用风险敞口
    \end{itemize}

    \textbf{D选项错误。}
    净额结算不是只在交易对手违约时才支付损失：
    \begin{itemize}
        \item 净额结算是正常的结算机制
        \item 不依赖于交易对手是否违约
        \item 这个描述不准确
    \end{itemize}
\end{explanation}

\checklist{净额结算允许公司将与特定交易对手的交易相互抵消}


\begin{question}
    Ondine Financial Inc. (Ondine) uses a variety of techniques to manage counterparty risk. It has entered into an interest rate swap with Scarbo Inc. (Scarbo). Currently, Ondine's position in the swap has a \$1.5 million mark-to-market value. Based on the information provided, which of the following credit risk mitigation techniques would be most advantageous to Ondine if Scarbo defaults?
\end{question}

\begin{choices}
    \item Close-out
    \item Collateralization
    \item Netting
    \item Walkaway
\end{choices}

\begin{correctanswer}
    D
\end{correctanswer}

\begin{explanation}
    \textbf{A选项错误。}
    终止（Close-out）在这种情况下对Ondine不利：
    \begin{itemize}
        \item Ondine的合约价值为负（-150万美元）
        \item 终止后需要支付净额
        \item 不能减少损失
    \end{itemize}

    \textbf{B选项错误。}
    抵押品（Collateralization）在这种情况下不相关：
    \begin{itemize}
        \item Ondine是净支付方
        \item 不需要提供抵押品
        \item 不能减少损失
    \end{itemize}

    \textbf{C选项错误。}
    净额结算（Netting）在这种情况下对Ondine不利：
    \begin{itemize}
        \item Ondine需要支付净额150万美元
        \item 不能减少损失
        \item 反而会增加损失
    \end{itemize}

    \textbf{D选项正确。}
    "撕毁"条款（Walkaway）在这种情况下对Ondine最有利：
    \begin{itemize}
        \item Ondine的合约价值为负
        \item 可以取消交易而不需要支付
        \item 可以最大限度地减少损失
    \end{itemize}
\end{explanation}

\checklist{当合约价值为负时，"撕毁"条款最能减少损失}



\begin{question}
    The following scenarios illustrate exposures for two trades with and without the impact of netting. Given the below expected positive exposures (EPE), what is the corresponding netting factor? Assume that the average correlation between exposures is 0.65.
\end{question}

\begin{choices}
    \item 0.65
    \item 0.75
    \item 0.85
    \item 1.35
\end{choices}

\begin{correctanswer}
    C
\end{correctanswer}

\begin{explanation}
    \textbf{A选项错误。}
    这个答案直接使用了相关系数，计算错误。净额因子需要考虑敞口数量和相关系数。

    \textbf{B选项错误。}
    这个答案计算错误。净额因子的计算需要考虑敞口数量和相关系数。

    \textbf{C选项正确。}
    净额因子的计算过程如下：
    \begin{itemize}
        \item 净额因子计算公式：\[ \text{netting factor} = \frac{\sqrt{n + n(n-1)\overline{\rho}}}{n} \]
        \item 其中：
        \begin{itemize}
            \item n = 敞口数量 = 5
            \item \(\overline{\rho}\) = 平均相关系数 = 0.65
        \end{itemize}
        \item 代入公式计算：
        \begin{itemize}
            \item \[ \text{Netting factor} = \frac{\sqrt{5 + 5(5-1) \times 0.65}}{5} \]
            \item \[ = \frac{\sqrt{5 + 5 \times 4 \times 0.65}}{5} \]
            \item \[ = \frac{\sqrt{5 + 13}}{5} \]
            \item \[ = \frac{\sqrt{18}}{5} \]
            \item \[ = \frac{4.24}{5} \]
            \item \[ = 0.85 \text{ 或 } 85\% \]
        \end{itemize}
    \end{itemize}

    \textbf{D选项错误。}
    这个答案计算错误。净额因子应该小于1，因为净额结算会减少风险敞口。
\end{explanation}

\checklist{净额因子等于根号下（敞口数量加上敞口数量乘以（敞口数量减1）乘以平均相关系数）除以敞口数量}


\begin{question}
    A firm has entered into a USD 25 million total return swap on the NASDAQ 100 Index as the index payer with ABC Corporation, which will pay 1-year LIBOR +3.0%. The contract will last 1 year, and cash flows will be exchanged annually. Suppose the NASDAQ 100 Index is currently at 3,200 and LIBOR is 1.50%. The firm conducts a stress test on this total return swap using the following scenario: NASDAQ 100 in 1 year: 4,000 LIBOR in 1 year: 0.75% For this scenario, what is the firm's net cash flow in year 1?
\end{question}

\begin{choices}
    \item A net cash outflow of USD 5.50 million
    \item A net cash outflow of USD 5.25 million
    \item A net cash outflow of USD 5.00 million
    \item A net cash outflow of USD 5.75 million
\end{choices}

\begin{correctanswer}
    B
\end{correctanswer}

\begin{explanation}
    \textbf{A选项错误。}
    这个答案计算错误。需要分别计算指数支付和LIBOR支付。

    \textbf{B选项正确。}
    计算过程如下：
    \begin{itemize}
        \item 指数支付：
        \begin{itemize}
            \item 指数上涨幅度 = (4,000/3,200) - 1 = 25%
            \item 指数支付金额 = 25% × \$25 million = \$6.25 million
        \end{itemize}
        \item LIBOR支付：
        \begin{itemize}
            \item 当前LIBOR = 1.50%
            \item 利差 = 3.0%
            \item 总利率 = 4.50%
            \item LIBOR支付金额 = 4.50% × \$25 million = \$1.00 million
        \end{itemize}
        \item 净现金流：
        \begin{itemize}
            \item 净现金流 = \$1.00 million - \$6.25 million = -\$5.25 million
            \item 负号表示现金流出
        \end{itemize}
    \end{itemize}

    \textbf{C选项错误。}
    这个答案计算错误。需要分别计算指数支付和LIBOR支付。

    \textbf{D选项错误。}
    这个答案计算错误。需要分别计算指数支付和LIBOR支付。
\end{explanation}

\checklist{总收益互换的净现金流等于LIBOR支付减去指数支付}




\newpage



\section{考点4:保证金（抵押品）和结算}
\setcounter{question}{0}


\begin{question}
    Bank A trades a 10-year interest rate swap with Bank B. Bank A is concerned about the 1-year exposure since Bank B has a poor credit rating. Which of the following measures will exacerbate Bank A's credit exposure to Bank B?
    I. Application of CSA
    II. Execute the swap as a reset swap wherein the swap will be marked to market every six months
    III. Application of break clause
    IV. Decrease the frequency of coupon payments
\end{question}

\begin{choices}
    \item I only
    \item IV only
    \item I, II, III and IV
    \item I, II and III
\end{choices}

\begin{correctanswer}
    B
\end{correctanswer}

\begin{explanation}
    \textbf{A选项错误。}
    信用支持协议（CSA）不会增加信用风险敞口：
    \begin{itemize}
        \item CSA要求提供抵押品
        \item 可以降低信用风险敞口
        \item 是风险缓释措施
    \end{itemize}

    \textbf{B选项正确。}
    减少息票支付频率会增加信用风险敞口：
    \begin{itemize}
        \item 支付间隔时间越长，市场利率变动的风险越大
        \item 增加了交易对手违约的风险
        \item 不利于风险控制
    \end{itemize}

    \textbf{C选项错误。}
    不是所有措施都会增加信用风险敞口：
    \begin{itemize}
        \item CSA、重置互换和终止条款都是风险缓释措施
        \item 只有减少支付频率会增加风险
        \item 这个选项错误地包含了所有措施
    \end{itemize}

    \textbf{D选项错误。}
    这个选项错误地包含了风险缓释措施：
    \begin{itemize}
        \item CSA、重置互换和终止条款都是风险缓释措施
        \item 这些措施不会增加信用风险敞口
        \item 只有减少支付频率会增加风险
    \end{itemize}
\end{explanation}

\checklist{减少息票支付频率会增加信用风险敞口}


\begin{question}
    Which of the following statements is least accurate regarding a credit support annex (CSA) and/or an ISDA Master Agreement?
\end{question}

\begin{choices}
    \item ISDA Master Agreements help standardize collateral management
    \item CSAs must define all collateralization parameters in order to work as intended
    \item Compared to the ISDA Master Agreement, CSAs were first to establish collateral standards
    \item CSAs are incorporated into an ISDA Master Agreement
\end{choices}

\begin{correctanswer}
    C
\end{correctanswer}

\begin{explanation}
    \textbf{A选项错误。}
    这个说法是正确的。ISDA主协议确实有助于标准化抵押品管理：
    \begin{itemize}
        \item 提供了标准化的法律框架
        \item 规定了抵押品管理的基本规则
        \item 促进了市场实践的统一
    \end{itemize}

    \textbf{B选项错误。}
    这个说法是正确的。信用支持附件（CSA）必须定义所有抵押品参数：
    \begin{itemize}
        \item 抵押品资格
        \item 利率支付
        \item 转移的时间和机制
        \item 抵押品计算
        \item 抵押品证券的折扣
        \item 抵押品替换
        \item 估值时间和方法
        \item 抵押品再利用
        \item 争议处理
    \end{itemize}

    \textbf{C选项正确。}
    这个说法是错误的。CSA不是最先建立抵押品标准的：
    \begin{itemize}
        \item ISDA主协议先于CSA建立
        \item CSA是作为ISDA主协议的补充而发展起来的
        \item 这个说法与历史发展不符
    \end{itemize}

    \textbf{D选项错误。}
    这个说法是正确的。CSA确实是ISDA主协议的组成部分：
    \begin{itemize}
        \item CSA是ISDA主协议的附件
        \item 用于规定抵押品管理的具体细节
        \item 可以选择单边或双边CSA
    \end{itemize}
\end{explanation}

\checklist{CSA不是最先建立抵押品标准的，而是作为ISDA主协议的补充而发展起来的}


\begin{question}
    When dealing with a hedge fund, a bank would most likely negotiate:
\end{question}

\begin{choices}
    \item A one-way agreement in the bank's favor given the bank's stronger credit rating
    \item A one-way agreement in the bank's favor agreeing to post collateral to the hedge fund
    \item A two-way agreement given the relatively small difference in credit quality between the two entities
    \item A two-way agreement where both parties agree to post collateral
\end{choices}

\begin{correctanswer}
    A
\end{correctanswer}

\begin{explanation}
    \textbf{A选项正确。}
    银行很可能会谈判一个对自己有利的单向协议：
    \begin{itemize}
        \item 银行的信用评级通常高于对冲基金
        \item 当交易双方信用质量差异较大时，适合签订单向协议
        \item 信用质量较差的一方（对冲基金）需要向信用质量较好的一方（银行）提供抵押品
    \end{itemize}

    \textbf{B选项错误。}
    这个选项描述的是不合理的安排：
    \begin{itemize}
        \item 银行信用质量更好，不应该向对冲基金提供抵押品
        \item 这种安排会增加银行的风险
        \item 不符合风险管理的原则
    \end{itemize}

    \textbf{C选项错误。}
    这个选项描述的情况不准确：
    \begin{itemize}
        \item 银行和对冲基金的信用质量通常有显著差异
        \item 不适合签订双向协议
        \item 这个选项低估了信用质量的差异
    \end{itemize}

    \textbf{D选项错误。}
    这个选项描述的是不合适的安排：
    \begin{itemize}
        \item 双向协议适用于信用质量相近的交易对手
        \item 银行和对冲基金的信用质量差异较大
        \item 这种安排不符合风险管理的最佳实践
    \end{itemize}
\end{explanation}

\checklist{当交易双方信用质量差异较大时，适合签订单向协议，信用质量较差的一方向信用质量较好的一方提供抵押品}


\begin{question}
    Leverage, Inc., an investment bank, has numerous credit default swaps with XYZ Corp. Leverage has established a break clause with XYZ Corp. to reduce risk. The break clause is trigger-based and may be exercised once the trigger is satisfied. The CEO of Leverage is concerned about a banker's paradox. Which of the following statements best describe the CEO's concern?
\end{question}

\begin{choices}
    \item To be effective, the break clause option should not be used too early
    \item The weak firm often recovers after the use of the break clause
    \item The break clause option is used too late, and the weak firm gets weaker
    \item The break clause option is used too early, and relations with the counterparty suffer
\end{choices}

\begin{correctanswer}
    C
\end{correctanswer}

\begin{explanation}
    \textbf{A选项错误。}
    这个说法与银行家悖论相反：
    \begin{itemize}
        \item 中断条款应该在交易对手信用质量显著下降之前使用
        \item 过早使用不是问题所在
        \item 这个选项误解了银行家悖论的本质
    \end{itemize}

    \textbf{B选项错误。}
    这个说法不准确：
    \begin{itemize}
        \item 交易对手在信用质量显著下降后通常不会恢复
        \item 这个选项低估了信用风险
        \item 不符合实际情况
    \end{itemize}

    \textbf{C选项正确。}
    这个说法准确描述了银行家悖论：
    \begin{itemize}
        \item 中断条款应该在交易对手信用质量显著下降之前使用
        \item 但机构通常为了避免损害与交易对手的关系而推迟使用
        \item 导致中断条款使用太晚，交易对手的信用质量进一步恶化
    \end{itemize}

    \textbf{D选项错误。}
    这个说法与银行家悖论相反：
    \begin{itemize}
        \item 过早使用不是主要问题
        \item 主要问题是使用太晚
        \item 这个选项误解了银行家悖论的本质
    \end{itemize}
\end{explanation}

\checklist{银行家悖论：中断条款应该在交易对手信用质量显著下降之前使用，但机构通常为了避免损害与交易对手的关系而推迟使用，导致使用太晚}

\begin{question}
    A financial institution has a two-way collateral support annex (CSA) with a counterparty covering a portfolio valued at JPY 500 million. The margining terms of the collateralized portfolio include a threshold of JPY 200 million, a minimum transfer amount of JPY 40 million, and a margin period of risk of 12 days. Which of the following is correct regarding the size of collateral in mitigating the counterparty risk of the portfolio?
\end{question}

\begin{choices}
    \item A lower threshold value is equivalent to a larger portion of exposure protected by collateral
    \item A shorter margin period of risk is equivalent to a smaller portion of exposure protected by collateral
    \item A lower independent amount is equivalent to a larger portion of exposure protected by collateral
    \item The protection from collateral specified in the CSA is uniform throughout the life of the exposure profile
\end{choices}

\begin{correctanswer}
    A
\end{correctanswer}

\begin{explanation}
    \textbf{A选项正确。}
    较低的门槛值意味着更大比例的敞口受到抵押品保护：
    \begin{itemize}
        \item 门槛值是未抵押敞口的金额
        \item 门槛值越低，需要提供抵押品的敞口比例越大
        \item 抵押品金额 = 敞口 - 门槛值
    \end{itemize}

    \textbf{B选项错误。}
    保证金风险期与抵押品金额无关：
    \begin{itemize}
        \item 保证金风险期是要求提供抵押品到收到抵押品之间的时间
        \item 不影响抵押品金额的计算
        \item 这个选项误解了保证金风险期的含义
    \end{itemize}

    \textbf{C选项错误。}
    较低的独立金额相当于较高的门槛值：
    \begin{itemize}
        \item 独立金额越低，门槛值越高
        \item 导致受抵押品保护的敞口比例越小
        \item 这个选项的因果关系错误
    \end{itemize}

    \textbf{D选项错误。}
    抵押品保护在整个敞口期间并不均匀：
    \begin{itemize}
        \item 在敞口较小的时候（开始和结束时）抵押品效果较小
        \item 抵押品保护效果随敞口大小变化
        \item 这个选项过于简化了抵押品保护的效果
    \end{itemize}
\end{explanation}

\checklist{门槛值越低，需要提供抵押品的敞口比例越大，抵押品金额等于敞口减去门槛值}


\begin{question}
    An underlying exposure with an effective annual price volatility of 7% is collateralized by a 10-year U.S. Treasury note with an effective price volatility of 9%. The correlation between the exposure on the U.S. Treasury note is zero. Changes in the value of the overall position (exposure plus collateral) are calculated for a 12-day horizon at a 95% confidence interval (assume a year of 250 days). Which of the following would one expect to observe from this analysis?
\end{question}

\begin{choices}
    \item The presence of collateral increases the current exposure and increase the volatility of the exposure between remargining periods
    \item The presence of collateral increases the current exposure, but decreases the volatility of the exposure between remargining periods
    \item The presence of collateral decreases the current exposure, but increases the volatility of the exposure between remargining periods
    \item The presence of collateral decreases the current exposure and decreases the volatility of the exposure between remargining periods
\end{choices}

\begin{correctanswer}
    C
\end{correctanswer}

\begin{explanation}
    \textbf{A选项错误。}
    这个选项描述的情况不正确：
    \begin{itemize}
        \item 抵押品会减少当前敞口
        \item 不会同时增加当前敞口和波动性
        \item 这个选项的因果关系错误
    \end{itemize}

    \textbf{B选项错误。}
    这个选项描述的情况不正确：
    \begin{itemize}
        \item 抵押品会减少当前敞口
        \item 不会减少波动性
        \item 这个选项的因果关系错误
    \end{itemize}

    \textbf{C选项正确。}
    计算过程如下：
    \begin{itemize}
        \item 抵押品的最坏情况变化：
        \begin{itemize}
            \item $-1.65 \times 9\% \times \sqrt{12/250} = -3.25\%$
        \end{itemize}
        \item 组合的总体波动性：
        \begin{itemize}
            \item $\sqrt{0.07^2 + 0.09^2} = 11.4\%$
        \end{itemize}
        \item 组合的最坏情况变化：
        \begin{itemize}
            \item $-1.65 \times 11.4\% \times \sqrt{12/250} = -4.12\%$
        \end{itemize}
        \item 因此，抵押品会减少当前敞口，但增加未来波动性
    \end{itemize}

    \textbf{D选项错误。}
    这个选项描述的情况不正确：
    \begin{itemize}
        \item 抵押品会减少当前敞口
        \item 但会增加波动性
        \item 这个选项的因果关系错误
    \end{itemize}
\end{explanation}

\checklist{抵押品会减少当前敞口，但会增加未来波动性，因为组合的总体波动性大于单个资产的波动性}


\begin{question}
    A financial institution has a two-way credit support annex with a counterparty covering a portfolio currently valued at USD 6,000,000. The margining terms of the collateralized portfolio include the following:
    · Threshold amount: USD 3,000,000
    · Minimum transfer amount: USD 500,000
    · Margin period of risk: 12 days
    Which of the following is correct regarding the impact of collateral on the counterparty risk of the portfolio?
\end{question}

\begin{choices}
    \item A shorter margin period of risk reduces the effectiveness of collateral in mitigating counterparty risk
    \item Imposing a minimum transfer amount can reduce counterparty risk
    \item Collateral has a large impact at the beginning but little impact at the end of the exposure profile
    \item A lower threshold value increases the effectiveness of collateral in mitigating counterparty risk
\end{choices}

\begin{correctanswer}
    D
\end{correctanswer}

\begin{explanation}
    \textbf{A选项错误。}
    这个说法不正确：
    \begin{itemize}
        \item 保证金风险期越短，抵押品效果越好
        \item 因为减少了抵押品延迟带来的风险
        \item 这个选项的因果关系错误
    \end{itemize}

    \textbf{B选项错误。}
    这个说法不正确：
    \begin{itemize}
        \item 最小转移金额是为了减少操作成本
        \item 不会减少对手方风险
        \item 这个选项误解了最小转移金额的作用
    \end{itemize}

    \textbf{C选项错误。}
    这个说法不正确：
    \begin{itemize}
        \item 抵押品在敞口较大时效果最显著
        \item 在敞口较小的时候（开始和结束时）效果较小
        \item 这个选项的因果关系错误
    \end{itemize}

    \textbf{D选项正确。}
    这个说法正确：
    \begin{itemize}
        \item 门槛值越低，需要提供抵押品的敞口比例越大
        \item 抵押品金额 = 敞口 - 门槛值
        \item 因此，较低的门槛值会增加抵押品在缓释对手方风险方面的效果
    \end{itemize}
\end{explanation}

\checklist{门槛值越低，需要提供抵押品的敞口比例越大，因此抵押品在缓释对手方风险方面的效果越好}



\newpage



\section{考点5:中央结算}
\setcounter{question}{0}

\begin{question}
    There are a number of challenges for clearing over-the-counter (OTC) derivative products through a centralized counterparty (CCP). Which of the following statements best summarizes the key challenges for central clearing of OTC derivative products?
\end{question}

\begin{choices}
    \item Jurisdictional fragmentation, increased counterparty risk, less transparency, and standardization
    \item Product complexity, illiquid products, presence of wrong-way risk, and lack of standardization
    \item Illiquid products, jurisdictional fragmentation, presence of wrong-way risk, and legal concerns
    \item Lack of standardization, increased counterparty risk, increased dependency risk, and less transparency
\end{choices}

\begin{correctanswer}
    B
\end{correctanswer}

\begin{explanation}
    \textbf{A选项错误。}
    这个选项描述不准确：
    \begin{itemize}
        \item 中央清算不会增加对手方风险
        \item 中央清算实际上增加了透明度
        \item 这个选项误解了中央清算的影响
    \end{itemize}

    \textbf{B选项正确。}
    这个选项准确描述了场外衍生品中央清算的主要挑战：
    \begin{itemize}
        \item 产品复杂性：场外衍生品结构复杂，难以估值
        \item 流动性低：某些产品缺乏流动性，难以在违约时平仓
        \item 错路风险：某些产品存在错路风险，增加了CCP的风险
        \item 缺乏标准化：场外衍生品需要标准化才能通过CCP清算
    \end{itemize}

    \textbf{C选项错误。}
    这个选项描述不完整：
    \begin{itemize}
        \item 虽然包含了部分正确因素
        \item 但忽略了产品复杂性和缺乏标准化这两个关键挑战
        \item 这个选项不够全面
    \end{itemize}

    \textbf{D选项错误。}
    这个选项描述不准确：
    \begin{itemize}
        \item 中央清算不会增加对手方风险
        \item 中央清算实际上增加了透明度
        \item 这个选项误解了中央清算的影响
    \end{itemize}
\end{explanation}

\checklist{场外衍生品中央清算的主要挑战包括产品复杂性、流动性低、错路风险和缺乏标准化}



\begin{question}
    Which of the following statements is an enhancement offered by the central counterparty (CCP) structure relative to the special purpose vehicle (SPV), the derivative product company (DPC), and the monoline insurance models?
\end{question}

\begin{choices}
    \item The CCP structure enables financial institutions to remove assets from their balance sheets
    \item The CCP structure enables counterparty risk to be outsourced, but in a non-diversified format
    \item The CCP structure spreads losses over a group of counterparties to minimize potential systemic risk
    \item The CCP structure enables a counterparty transaction originator to fail and not affect the other member firms
\end{choices}

\begin{correctanswer}
    C
\end{correctanswer}

\begin{explanation}
    \textbf{A选项错误。}
    这个说法不正确：
    \begin{itemize}
        \item 从资产负债表移除资产是SPV和DPC的功能
        \item CCP不提供这种功能
        \item 这个选项混淆了不同机构的功能
    \end{itemize}

    \textbf{B选项错误。}
    这个说法不正确：
    \begin{itemize}
        \item 以非分散方式外包对手方风险是单一险种保险公司的功能
        \item CCP提供的是分散化的风险转移
        \item 这个选项混淆了不同机构的功能
    \end{itemize}

    \textbf{C选项正确。}
    这个说法正确：
    \begin{itemize}
        \item CCP通过抵押品和损失共担机制将损失分散在一组交易对手中
        \item 这种机制降低了潜在的系统性风险
        \item 这是CCP相对于其他模式的主要优势
    \end{itemize}

    \textbf{D选项错误。}
    这个说法不正确：
    \begin{itemize}
        \item 在CCP结构中，如果交易发起方违约，所有成员公司都会受到影响
        \item 通过损失瀑布机制（loss waterfall）分配损失
        \item 这个选项误解了CCP的运作机制
    \end{itemize}
\end{explanation}

\checklist{CCP通过抵押品和损失共担机制将损失分散在一组交易对手中，降低了潜在的系统性风险}


\begin{question}
    Which of the following statements is not an improvement that centrally cleared markets offer relative to bilateral markets?
\end{question}

\begin{choices}
    \item Centrally cleared markets improve the counterparty risk picture by replacing the original counterparty with a series of counterparties
    \item Centrally cleared markets remain market-neutral by netting trades
    \item Centrally cleared markets formalize the default work-out process by utilizing a loss waterfall structure
    \item Centrally cleared markets offer more flexibility in contract selection because of their collateral collecting process
\end{choices}

\begin{correctanswer}
    D
\end{correctanswer}

\begin{explanation}
    \textbf{A选项错误。}
    这个说法是正确的：
    \begin{itemize}
        \item 中央清算市场通过将原始对手方替换为一系列对手方来改善对手方风险状况
        \item 这种多边结构降低了单个对手方违约的影响
        \item 这是中央清算市场的一个优势
    \end{itemize}

    \textbf{B选项错误。}
    这个说法是正确的：
    \begin{itemize}
        \item 中央清算市场通过净额结算保持市场中性
        \item 减少了交易对手之间的支付金额
        \item 这是中央清算市场的一个优势
    \end{itemize}

    \textbf{C选项错误。}
    这个说法是正确的：
    \begin{itemize}
        \item 中央清算市场通过损失瀑布结构使违约处理过程规范化
        \item 提供了清晰的损失分配机制
        \item 这是中央清算市场的一个优势
    \end{itemize}

    \textbf{D选项正确。}
    这个说法是错误的：
    \begin{itemize}
        \item 双边市场允许任何类型的定制金融合约和定制抵押品
        \item 中央清算市场要求合约标准化，抵押品规则固定且不可协商
        \item 因此，中央清算市场实际上降低了合约选择的灵活性
    \end{itemize}
\end{explanation}

\checklist{中央清算市场要求合约标准化，抵押品规则固定且不可协商，因此降低了合约选择的灵活性}


\begin{question}
    Which of the following statements is an enhancement offered by the central counterparty (CCP) structure relative to the special purpose vehicle (SPV), the derivative product company (DPC), and the monoline insurance models?
\end{question}

\begin{choices}
    \item The CCP structure enables financial institutions to remove assets from their balance sheets
    \item The CCP structure enables counterparty risk to be outsourced, but in a non-diversified format
    \item The CCP structure spreads losses over a group of counterparties to minimize potential systemic risk
    \item The CCP structure enables a counterparty transaction originator to fail and not affect the other member firms
\end{choices}

\begin{correctanswer}
    C
\end{correctanswer}

\begin{explanation}
    \textbf{A选项错误。}
    这个说法不正确：
    \begin{itemize}
        \item 从资产负债表移除资产是SPV和DPC的功能
        \item CCP不提供这种功能
        \item 这个选项混淆了不同机构的功能
    \end{itemize}

    \textbf{B选项错误。}
    这个说法不正确：
    \begin{itemize}
        \item 以非分散方式外包对手方风险是单一险种保险公司的功能
        \item CCP提供的是分散化的风险转移
        \item 这个选项混淆了不同机构的功能
    \end{itemize}

    \textbf{C选项正确。}
    这个说法正确：
    \begin{itemize}
        \item CCP通过抵押品和损失共担机制将损失分散在一组交易对手中
        \item 这种机制降低了潜在的系统性风险
        \item 这是CCP相对于其他模式的主要优势
    \end{itemize}

    \textbf{D选项错误。}
    这个说法不正确：
    \begin{itemize}
        \item 在CCP结构中，如果交易发起方违约，所有成员公司都会受到影响
        \item 通过损失瀑布机制（loss waterfall）分配损失
        \item 这个选项误解了CCP的运作机制
    \end{itemize}
\end{explanation}

\checklist{CCP通过抵押品和损失共担机制将损失分散在一组交易对手中，降低了潜在的系统性风险}




\newpage



\section{考点6:未来值和敞口}
\setcounter{question}{0}

\begin{question}
    Each of the following is true except?
\end{question}

\begin{choices}
    \item At any point in time, effective EE cannot be less than EE
    \item (effective) EPE is average (effective) EE over time
    \item Effective EPE cannot be less than EPE
    \item For each point in time, there is a different maximum PFE such that maximum PFE does not represent a single value
\end{choices}

\begin{correctanswer}
    D
\end{correctanswer}

\begin{explanation}
    \textbf{A选项错误。}
    这个说法是正确的：
    \begin{itemize}
        \item 有效期望敞口（effective EE）在任何时点都不能小于期望敞口（EE）
        \item 因为有效期望敞口考虑了净额结算的影响
        \item 净额结算会减少敞口，但不会增加敞口
    \end{itemize}

    \textbf{B选项错误。}
    这个说法是正确的：
    \begin{itemize}
        \item 有效期望正敞口（effective EPE）是有效期望敞口（effective EE）在时间上的平均值
        \item 反映了整个期间的平均风险敞口
        \item 这是期望正敞口的基本定义
    \end{itemize}

    \textbf{C选项错误。}
    这个说法是正确的：
    \begin{itemize}
        \item 有效期望正敞口（effective EPE）不能小于期望正敞口（EPE）
        \item 因为有效期望正敞口考虑了净额结算的影响
        \item 净额结算会减少敞口，但不会增加敞口
    \end{itemize}

    \textbf{D选项正确。}
    这个说法是错误的：
    \begin{itemize}
        \item 最大潜在未来敞口（maximum PFE）是一个单一值
        \item 它代表给定时间区间内最高的PFE值
        \item 这个选项误解了maximum PFE的定义
    \end{itemize}
\end{explanation}

\checklist{最大潜在未来敞口（maximum PFE）是一个单一值，代表给定时间区间内最高的PFE值}

\begin{question}
    Which metric for credit exposure is represented by the ? in the following graph?
    \begin{figure}[htbp]
        \centering
        \includegraphics[width=0.4\textwidth]{images/1.png}
        \caption{}
    \end{figure}
\end{question}

\begin{choices}
    \item Expected positive exposure (EPE)
    \item Potential future exposure (PFE)
    \item Effective EE
    \item Effective EPE
\end{choices}

\begin{correctanswer}
    A
\end{correctanswer}

\begin{explanation}
    \textbf{A选项正确。}
    期望正敞口（EPE）是正确的答案：
    \begin{itemize}
        \item EPE是期望敞口（EE）在时间上的平均值
        \item 图中问号表示的是EE曲线下方的平均面积
        \item 这是EPE的图形表示方式
    \end{itemize}

    \textbf{B选项错误。}
    潜在未来敞口（PFE）不是正确答案：
    \begin{itemize}
        \item PFE是在给定置信水平下的极端敞口
        \item 图中问号不是表示PFE
        \item 这个选项误解了图形的含义
    \end{itemize}

    \textbf{C选项错误。}
    有效期望敞口（Effective EE）不是正确答案：
    \begin{itemize}
        \item Effective EE假设信用敞口是非递减的
        \item 图中问号不是表示Effective EE
        \item 这个选项误解了图形的含义
    \end{itemize}

    \textbf{D选项错误。}
    有效期望正敞口（Effective EPE）不是正确答案：
    \begin{itemize}
        \item Effective EPE是Effective EE的均值
        \item 图中问号不是表示Effective EPE
        \item 这个选项误解了图形的含义
    \end{itemize}
\end{explanation}

\checklist{期望正敞口（EPE）是期望敞口（EE）在时间上的平均值，在图中表示为EE曲线下方的平均面积}


\begin{question}
    When short an OTC put option on ABC stock, how does the credit exposure change if the stock price decreases?
\end{question}

\begin{choices}
    \item Increases
    \item Decreases
    \item Doesn't vary with underlying stock price
    \item There is no credit exposure on short options
\end{choices}

\begin{correctanswer}
    D
\end{correctanswer}

\begin{explanation}
    \textbf{A选项错误。}
    这个说法不正确：
    \begin{itemize}
        \item 期权空头头寸在期初已经收到保费
        \item 不会因为标的股票价格下跌而产生信用风险敞口
        \item 这个选项误解了期权空头头寸的信用风险特征
    \end{itemize}

    \textbf{B选项错误。}
    这个说法不正确：
    \begin{itemize}
        \item 期权空头头寸在期初已经收到保费
        \item 不会因为标的股票价格下跌而减少信用风险敞口
        \item 这个选项误解了期权空头头寸的信用风险特征
    \end{itemize}

    \textbf{C选项错误。}
    这个说法不正确：
    \begin{itemize}
        \item 期权空头头寸在期初已经收到保费
        \item 不会因为标的股票价格变化而产生信用风险敞口
        \item 这个选项误解了期权空头头寸的信用风险特征
    \end{itemize}

    \textbf{D选项正确。}
    这个说法正确：
    \begin{itemize}
        \item 期权空头头寸在期初已经收到保费
        \item 不会产生信用风险敞口
        \item 因为期权卖方已经收到了全部收益
    \end{itemize}
\end{explanation}

\checklist{期权空头头寸在期初已经收到保费，因此不会产生信用风险敞口}


\begin{question}
    BNP Paribas has just entered into a plain-vanilla interest-rate swap as a pay-fixed counterparty. Credit Agricole is the receive-fixed counterparty in the same swap. The forward spot curve is upward-sloping. If LIBOR starts trending down and the forward spot curve flattens, the credit risk from the swap will:
\end{question}

\begin{choices}
    \item Increase only for BNP Paribas
    \item Increase only for Credit Agricole
    \item Decrease for both BNP Paribas and Credit Agricole
    \item Increase for both BNP Paribas and Credit Agricole
\end{choices}

\begin{correctanswer}
    B
\end{correctanswer}

\begin{explanation}
    \textbf{A选项错误。}
    这个说法不正确：
    \begin{itemize}
        \item BNP Paribas是支付固定利率的一方
        \item 当LIBOR下降时，BNP Paribas的信用风险会减少
        \item 因为BNP Paribas需要支付的金额相对增加
    \end{itemize}

    \textbf{B选项正确。}
    这个说法正确：
    \begin{itemize}
        \item Credit Agricole是收取固定利率的一方
        \item 当LIBOR下降时，Credit Agricole支付的浮动利率金额减少
        \item 导致Credit Agricole的信用风险敞口增加
    \end{itemize}

    \textbf{C选项错误。}
    这个说法不正确：
    \begin{itemize}
        \item 当LIBOR下降时，BNP Paribas的信用风险会减少
        \item 但Credit Agricole的信用风险会增加
        \item 这个选项误解了利率变化对双方信用风险的影响
    \end{itemize}

    \textbf{D选项错误。}
    这个说法不正确：
    \begin{itemize}
        \item 当LIBOR下降时，BNP Paribas的信用风险会减少
        \item 只有Credit Agricole的信用风险会增加
        \item 这个选项误解了利率变化对双方信用风险的影响
    \end{itemize}
\end{explanation}

\checklist{在利率互换中，当LIBOR下降时，收取固定利率的一方的信用风险会增加，因为其支付的浮动利率金额减少}


\begin{question}
    The potential future credit exposure profile peaks at maturity for which of the following instruments:
    I. FX forwards
    II. Interest rate swaps
    III. Cross currency swaps with final exchange
\end{question}

\begin{choices}
    \item I only
    \item II only
    \item I and III
    \item I, II and III
\end{choices}

\begin{correctanswer}
    C
\end{correctanswer}

\begin{explanation}
    \textbf{A选项错误。}
    这个选项不完整：
    \begin{itemize}
        \item 外汇远期（FX forwards）的信用风险敞口在到期时达到峰值
        \item 但跨币种互换（cross currency swaps）也有相同特征
        \item 这个选项遗漏了跨币种互换
    \end{itemize}

    \textbf{B选项错误。}
    这个选项不正确：
    \begin{itemize}
        \item 利率互换（interest rate swaps）的信用风险敞口在到期前达到峰值
        \item 因为不涉及本金的交换
        \item 这个选项误解了利率互换的风险特征
    \end{itemize}

    \textbf{C选项正确。}
    这个选项正确：
    \begin{itemize}
        \item 外汇远期和跨币种互换在到期时需要交换本金
        \item 因此信用风险敞口在到期时达到峰值
        \item 利率互换不涉及本金交换，风险敞口在到期前达到峰值
    \end{itemize}

    \textbf{D选项错误。}
    这个选项不正确：
    \begin{itemize}
        \item 利率互换的信用风险敞口在到期前达到峰值
        \item 不应该包含在答案中
        \item 这个选项错误地包含了利率互换
    \end{itemize}
\end{explanation}

\checklist{外汇远期和跨币种互换的信用风险敞口在到期时达到峰值，因为需要交换本金；而利率互换的信用风险敞口在到期前达到峰值，因为不涉及本金交换}


\begin{question}
    Which of the following security types will most likely result in a peaked shape for the exposure profile represented by potential future exposure (PFE)?
\end{question}

\begin{choices}
    \item Long option position
    \item Foreign exchange product
    \item 10-year loan with a floating rate payment
    \item Swap
\end{choices}

\begin{correctanswer}
    D
\end{correctanswer}

\begin{explanation}
    \textbf{A选项错误。}
    这个选项不正确：
    \begin{itemize}
        \item 期权多头头寸的信用风险敞口不会呈现峰值形状
        \item 期权多头头寸的信用风险敞口通常随时间递减
        \item 因为期权价值随时间衰减
    \end{itemize}

    \textbf{B选项错误。}
    这个选项不正确：
    \begin{itemize}
        \item 外汇产品的信用风险敞口在到期时达到峰值
        \item 不会呈现中间峰值的形状
        \item 因为需要在到期时交换本金
    \end{itemize}

    \textbf{C选项错误。}
    这个选项不正确：
    \begin{itemize}
        \item 浮动利率贷款的信用风险敞口通常随时间递减
        \item 不会呈现峰值形状
        \item 因为本金余额随时间减少
    \end{itemize}

    \textbf{D选项正确。}
    这个选项正确：
    \begin{itemize}
        \item 互换的信用风险敞口呈现峰值形状
        \item 前期受到现金流不确定性的影响，敞口上升
        \item 后期随着交割完成，现金流减少，敞口下降
    \end{itemize}
\end{explanation}

\checklist{互换的信用风险敞口呈现峰值形状，前期因现金流不确定性而上升，后期因交割完成而下降}


\begin{question}
    Which of the following statements regarding exposure measures is not accurate?
\end{question}

\begin{choices}
    \item Assuming no recovery in bankruptcy, current exposure is the greater of zero and the market value of a transaction or portfolio of transactions that the investor would be lost when a counterparty defaulted
    \item Potential Future Exposure measures the distribution of exposures of low percentile
    \item Expected exposure measures the mean distribution of exposures at a given future date prior to the maturity of the longest maturity exposure in the netting group
    \item The weights of expected positive exposure could represent the proportion that an individual expected exposures of the entire time interval
\end{choices}

\begin{correctanswer}
    B
\end{correctanswer}

\begin{explanation}
    \textbf{A选项错误。}
    这个说法是正确的：
    \begin{itemize}
        \item 当前敞口是在交易对手违约时投资者可能损失的市场价值
        \item 假设破产时没有回收
        \item 取零和市场价值的较大值
    \end{itemize}

    \textbf{B选项正确。}
    这个说法是错误的：
    \begin{itemize}
        \item 潜在未来敞口（PFE）衡量的是高百分位数的敞口分布
        \item 而不是低百分位数的敞口分布
        \item 这个选项误解了PFE的定义
    \end{itemize}

    \textbf{C选项错误。}
    这个说法是正确的：
    \begin{itemize}
        \item 期望敞口衡量的是在给定未来日期的敞口均值分布
        \item 时间范围在净额组合中最长期限之前
        \item 这个选项准确描述了期望敞口的定义
    \end{itemize}

    \textbf{D选项错误。}
    这个说法是正确的：
    \begin{itemize}
        \item 期望正敞口的权重可以代表单个期望敞口在整个时间区间中的比例
        \item 这是期望正敞口计算的基本原理
        \item 这个选项准确描述了期望正敞口的权重含义
    \end{itemize}
\end{explanation}

\checklist{潜在未来敞口（PFE）衡量的是高百分位数的敞口分布，而不是低百分位数的敞口分布}

\begin{question}
    Each of the following is true except?
\end{question}

\begin{choices}
    \item At any point in time, effective EE cannot be less than EE
    \item (effective) EPE is average (effective) EE over time
    \item Effective EPE cannot be less than EPE
    \item For each point in time, there is a different maximum PFE such that maximum PFE does not represent a single value
\end{choices}

\begin{correctanswer}
    D
\end{correctanswer}

\begin{explanation}
    \textbf{A选项错误。}
    这个说法是正确的：
    \begin{itemize}
        \item 有效期望敞口（effective EE）在任何时点都不能小于期望敞口（EE）
        \item 因为有效期望敞口考虑了净额结算的影响
        \item 净额结算会减少敞口，但不会增加敞口
    \end{itemize}

    \textbf{B选项错误。}
    这个说法是正确的：
    \begin{itemize}
        \item 有效期望正敞口（effective EPE）是有效期望敞口（effective EE）在时间上的平均值
        \item 反映了整个期间的平均风险敞口
        \item 这是期望正敞口的基本定义
    \end{itemize}

    \textbf{C选项错误。}
    这个说法是正确的：
    \begin{itemize}
        \item 有效期望正敞口（effective EPE）不能小于期望正敞口（EPE）
        \item 因为有效期望正敞口考虑了净额结算的影响
        \item 净额结算会减少敞口，但不会增加敞口
    \end{itemize}

    \textbf{D选项正确。}
    这个说法是错误的：
    \begin{itemize}
        \item 最大潜在未来敞口（maximum PFE）是一个单一值
        \item 它代表给定时间区间内最高的PFE值
        \item 这个选项误解了maximum PFE的定义
    \end{itemize}
\end{explanation}

\checklist{最大潜在未来敞口（maximum PFE）是一个单一值，代表给定时间区间内最高的PFE值}


\begin{question}
    Which of the following security types will most likely result in a peaked shape for the exposure profile represented by potential future exposure (PFE)?
\end{question}

\begin{choices}
    \item Long option position
    \item Foreign exchange product
    \item 10-year loan with a floating rate payment
    \item Swap
\end{choices}

\begin{correctanswer}
    D
\end{correctanswer}

\begin{explanation}
    \textbf{A选项错误。}
    这个选项不正确：
    \begin{itemize}
        \item 期权多头头寸的信用风险敞口不会呈现峰值形状
        \item 期权多头头寸的信用风险敞口通常随时间递减
        \item 因为期权价值随时间衰减
    \end{itemize}

    \textbf{B选项错误。}
    这个选项不正确：
    \begin{itemize}
        \item 外汇产品的信用风险敞口在到期时达到峰值
        \item 不会呈现中间峰值的形状
        \item 因为需要在到期时交换本金
    \end{itemize}

    \textbf{C选项错误。}
    这个选项不正确：
    \begin{itemize}
        \item 浮动利率贷款的信用风险敞口通常随时间递减
        \item 不会呈现峰值形状
        \item 因为本金余额随时间减少
    \end{itemize}

    \textbf{D选项正确。}
    这个选项正确：
    \begin{itemize}
        \item 互换的信用风险敞口呈现峰值形状
        \item 前期受到现金流不确定性的影响，敞口上升
        \item 后期随着交割完成，现金流减少，敞口下降
    \end{itemize}
\end{explanation}

\checklist{互换的信用风险敞口呈现峰值形状，前期因现金流不确定性而上升，后期因交割完成而下降}


\begin{question}
    BNP Paribas has just entered into a plain-vanilla interest-rate swap as a pay-fixed counterparty. Credit Agricole is the receive-fixed counterparty in the same swap. The forward spot curve is upward-sloping. If LIBOR starts trending down and the forward spot curve flattens, the credit risk from the swap will:
\end{question}

\begin{choices}
    \item Increase only for BNP Paribas
    \item Increase only for Credit Agricole
    \item Decrease for both BNP Paribas and Credit Agricole
    \item Increase for both BNP Paribas and Credit Agricole
\end{choices}

\begin{correctanswer}
    B
\end{correctanswer}

\begin{explanation}
    \textbf{A选项错误。}
    这个说法不正确：
    \begin{itemize}
        \item BNP Paribas是支付固定利率的一方
        \item 当LIBOR下降时，BNP Paribas的信用风险会减少
        \item 因为BNP Paribas需要支付的金额相对增加
    \end{itemize}

    \textbf{B选项正确。}
    这个说法正确：
    \begin{itemize}
        \item Credit Agricole是收取固定利率的一方
        \item 当LIBOR下降时，Credit Agricole支付的浮动利率金额减少
        \item 导致Credit Agricole的信用风险敞口增加
    \end{itemize}

    \textbf{C选项错误。}
    这个说法不正确：
    \begin{itemize}
        \item 当LIBOR下降时，BNP Paribas的信用风险会减少
        \item 但Credit Agricole的信用风险会增加
        \item 这个选项误解了利率变化对双方信用风险的影响
    \end{itemize}

    \textbf{D选项错误。}
    这个说法不正确：
    \begin{itemize}
        \item 当LIBOR下降时，BNP Paribas的信用风险会减少
        \item 只有Credit Agricole的信用风险会增加
        \item 这个选项误解了利率变化对双方信用风险的影响
    \end{itemize}
\end{explanation}

\checklist{在利率互换中，当LIBOR下降时，收取固定利率的一方的信用风险会增加，因为其支付的浮动利率金额减少}





\newpage



\section{考点7:CVA}
\setcounter{question}{0}
\begin{question}
    Which of the following statements regarding counterparty credit risk are correct?
    \begin{enumerate}
        \item Expected positive exposure is the highest expected exposure over a specified interval
        \item Wrong-way exposures are positively correlated with the counterparty's credit quality
        \item Credit triggers are early settlement agreements that require counterparties to settle and terminate trades if the credit rating of a party falls below a specified level
        \item Right-way exposures are negatively correlated with the counterparty's credit quality
        \item Cross-product netting is a provision that allows counterparties to net payments across different products
        \item Collateral agreements require that specified amounts of liabilities be transferred to counterparty if exposures exceed a specified threshold
    \end{enumerate}
\end{question}

\begin{choices}
    \item III and V only
    \item I, III, and V
    \item I, II, and IV
    \item III, V, and VI
\end{choices}

\begin{correctanswer}
    A
\end{correctanswer}

\begin{explanation}
    \textbf{A选项正确。}
    这个选项正确：
    \begin{itemize}
        \item 信用触发条款（Credit triggers）要求当一方信用评级低于特定水平时，交易对手需要结算并终止交易
        \item 跨产品净额结算（Cross-product netting）允许交易对手在不同产品之间进行净额结算
        \item 其他说法都是错误的
    \end{itemize}

    \textbf{B选项错误。}
    这个选项不正确：
    \begin{itemize}
        \item 期望正敞口（EPE）是特定区间内的平均期望敞口，而不是最高期望敞口
        \item 其他说法正确
        \item 这个选项包含了错误的说法
    \end{itemize}

    \textbf{C选项错误。}
    这个选项不正确：
    \begin{itemize}
        \item 错路风险（Wrong-way risk）与交易对手信用质量呈负相关
        \item 对路风险（Right-way risk）与交易对手信用质量呈正相关
        \item 这个选项误解了错路风险和对路风险的定义
    \end{itemize}

    \textbf{D选项错误。}
    这个选项不正确：
    \begin{itemize}
        \item 抵押协议要求转移资产而不是负债
        \item 其他说法正确
        \item 这个选项包含了错误的说法
    \end{itemize}
\end{explanation}

\checklist{信用触发条款要求当一方信用评级低于特定水平时结算并终止交易；跨产品净额结算允许在不同产品之间进行净额结算；错路风险与交易对手信用质量呈负相关，对路风险与交易对手信用质量呈正相关}


\begin{question}
    Regarding the impact of changes in the credit spread and recovery rate assumptions on the CVA, which of the following statements is true?
\end{question}

\begin{choices}
    \item A decrease in the credit spread will most often increase the CVA
    \item For an upward-sloping curve, the CVA will be higher compared to a downward-sloping curve
    \item Increasing the recovery rate will reduce the CVA
    \item If the actual recovery rate is higher than the settled recovery rate, the CVA will most likely be higher compared to a situation where both recovery assumptions are the same for both rates
\end{choices}

\begin{correctanswer}
    C
\end{correctanswer}

\begin{explanation}
    \textbf{A选项错误。}
    这个说法不正确：
    \begin{itemize}
        \item 信用利差下降通常会降低CVA
        \item 因为违约概率降低
        \item 这个选项误解了信用利差对CVA的影响
    \end{itemize}

    \textbf{B选项错误。}
    这个说法不正确：
    \begin{itemize}
        \item 向上倾斜的信用利差曲线会导致较低的CVA
        \item 相对于向下倾斜的曲线
        \item 这个选项误解了信用利差曲线形状对CVA的影响
    \end{itemize}

    \textbf{C选项正确。}
    这个说法正确：
    \begin{itemize}
        \item 提高回收率会增加隐含违约概率
        \item 但会降低最终的CVA
        \item 因为违约损失减少
    \end{itemize}

    \textbf{D选项错误。}
    这个说法不正确：
    \begin{itemize}
        \item 实际回收率高于设定回收率会导致较低的CVA
        \item 相对于两者相同的情况
        \item 这个选项误解了回收率差异对CVA的影响
    \end{itemize}
\end{explanation}

\checklist{提高回收率会增加隐含违约概率但会降低CVA，因为违约损失减少；向上倾斜的信用利差曲线会导致较低的CVA；实际回收率高于设定回收率会导致较低的CVA}


\begin{question}
    With respect to the CVA calculation, which of the following statements is correct when a risk manager wishes to understand which trades have the greatest impact on a counterparty's CVA? The manager would use:
\end{question}

\begin{choices}
    \item Incremental CVA because it accounts for the change in CVA once the new trade is priced, accounting for netting
    \item Marginal CVA because he could break down netted trades into trade level contributions
    \item Incremental CVA because he could break down netted trades into trade level contributions
    \item Marginal CVA because it accounts for the change in CVA once the new trade is priced, accounting for netting
\end{choices}

\begin{correctanswer}
    B
\end{correctanswer}

\begin{explanation}
    \textbf{A选项错误。}
    这个说法不正确：
    \begin{itemize}
        \item 增量CVA（Incremental CVA）用于评估新交易对CVA的影响
        \item 不能用于分析现有交易对CVA的贡献
        \item 这个选项误解了增量CVA的用途
    \end{itemize}

    \textbf{B选项正确。}
    这个说法正确：
    \begin{itemize}
        \item 边际CVA（Marginal CVA）可以将净额结算后的交易分解为单个交易的贡献
        \item 帮助风险经理了解哪些交易对CVA影响最大
        \item 这是边际CVA的主要用途
    \end{itemize}

    \textbf{C选项错误。}
    这个说法不正确：
    \begin{itemize}
        \item 增量CVA不能将净额结算后的交易分解为单个交易的贡献
        \item 增量CVA只用于评估新交易的影响
        \item 这个选项误解了增量CVA的功能
    \end{itemize}

    \textbf{D选项错误。}
    这个说法不正确：
    \begin{itemize}
        \item 边际CVA不是用于评估新交易对CVA的影响
        \item 边际CVA用于分析现有交易的贡献
        \item 这个选项误解了边际CVA的用途
    \end{itemize}
\end{explanation}

\checklist{边际CVA可以将净额结算后的交易分解为单个交易的贡献，帮助风险经理了解哪些交易对CVA影响最大；而增量CVA用于评估新交易对CVA的影响}


\begin{question}
    Regarding the impact of changes in the credit spread and recovery rate assumptions on the CVA, which of the following statements is true?
\end{question}

\begin{choices}
    \item A decrease in the credit spread will most often increase the CVA
    \item For an upward-sloping curve, the CVA will be higher compared to a downward-sloping curve
    \item Increasing the recovery rate will reduce the CVA
    \item If the actual recovery rate is higher than the settled recovery rate, the CVA will most likely be higher compared to a situation where both recovery assumptions are the same for both rates
\end{choices}

\begin{correctanswer}
    C
\end{correctanswer}

\begin{explanation}
    \textbf{A选项错误。}
    这个说法不正确：
    \begin{itemize}
        \item 信用利差下降通常会降低CVA
        \item 因为违约概率降低
        \item 这个选项误解了信用利差对CVA的影响
    \end{itemize}

    \textbf{B选项错误。}
    这个说法不正确：
    \begin{itemize}
        \item 向上倾斜的信用利差曲线会导致较低的CVA
        \item 相对于向下倾斜的曲线
        \item 这个选项误解了信用利差曲线形状对CVA的影响
    \end{itemize}

    \textbf{C选项正确。}
    这个说法正确：
    \begin{itemize}
        \item 提高回收率会增加隐含违约概率
        \item 但会降低最终的CVA
        \item 因为违约损失减少
    \end{itemize}

    \textbf{D选项错误。}
    这个说法不正确：
    \begin{itemize}
        \item 实际回收率高于设定回收率会导致较低的CVA
        \item 相对于两者相同的情况
        \item 这个选项误解了回收率差异对CVA的影响
    \end{itemize}
\end{explanation}

\checklist{回收率提高降低CVA，信用利差曲线向上倾斜降低CVA}

\begin{question}
    With respect to the CVA calculation, which of the following statements is correct when a risk manager wishes to understand which trades have the greatest impact on a counterparty's CVA? The manager would use:
\end{question}

\begin{choices}
    \item Incremental CVA because it accounts for the change in CVA once the new trade is priced, accounting for netting
    \item Marginal CVA because he could break down netted trades into trade level contributions
    \item Incremental CVA because he could break down netted trades into trade level contributions
    \item Marginal CVA because it accounts for the change in CVA once the new trade is priced, accounting for netting
\end{choices}

\begin{correctanswer}
    B
\end{correctanswer}

\begin{explanation}
    \textbf{A选项错误。}
    这个说法不正确：
    \begin{itemize}
        \item 增量CVA（Incremental CVA）用于评估新交易对CVA的影响
        \item 不能用于分析现有交易对CVA的贡献
        \item 这个选项误解了增量CVA的用途
    \end{itemize}

    \textbf{B选项正确。}
    这个说法正确：
    \begin{itemize}
        \item 边际CVA（Marginal CVA）可以将净额结算后的交易分解为单个交易的贡献
        \item 帮助风险经理了解哪些交易对CVA影响最大
        \item 这是边际CVA的主要用途
    \end{itemize}

    \textbf{C选项错误。}
    这个说法不正确：
    \begin{itemize}
        \item 增量CVA不能将净额结算后的交易分解为单个交易的贡献
        \item 增量CVA只用于评估新交易的影响
        \item 这个选项误解了增量CVA的功能
    \end{itemize}

    \textbf{D选项错误。}
    这个说法不正确：
    \begin{itemize}
        \item 边际CVA不是用于评估新交易对CVA的影响
        \item 边际CVA用于分析现有交易的贡献
        \item 这个选项误解了边际CVA的用途
    \end{itemize}
\end{explanation}

\checklist{边际CVA用于分析单个交易对CVA的贡献}


\begin{question}
    An investor said: "The formula for BCVA is similar to the formula for CVA, except that the BCVA formula incorporates the probability of the counterparty's survival and uses EPE." Which of the following statements is accurate?
\end{question}

\begin{choices}
    \item The statement is correct with regard to EPE only
    \item The statement is correct with regard to probability of survival only
    \item The statement is incorrect with regard to both EPE and probability of survival
    \item The statement is correct with regard to both EPE and probability of survival
\end{choices}

\begin{correctanswer}
    B
\end{correctanswer}

\begin{explanation}
    \textbf{A选项错误。}
    这个说法不正确：
    \begin{itemize}
        \item BCVA公式使用负期望敞口（NEE）而不是EPE
        \item 这个选项误解了BCVA公式中使用的敞口指标
        \item 这个选项部分正确但不够完整
    \end{itemize}

    \textbf{B选项正确。}
    这个说法正确：
    \begin{itemize}
        \item BCVA公式确实需要考虑交易对手的生存概率
        \item 这是BCVA与CVA的主要区别之一
        \item 这个选项准确描述了BCVA公式的特征
    \end{itemize}

    \textbf{C选项错误。}
    这个说法不正确：
    \begin{itemize}
        \item 关于生存概率的说法是正确的
        \item 只有关于EPE的说法是错误的
        \item 这个选项完全否定了正确的部分
    \end{itemize}

    \textbf{D选项错误。}
    这个说法不正确：
    \begin{itemize}
        \item 关于EPE的说法是错误的
        \item BCVA使用NEE而不是EPE
        \item 这个选项错误地认为两个说法都正确
    \end{itemize}
\end{explanation}

\checklist{BCVA使用NEE而非EPE，但需考虑对手方生存概率}


\begin{question}
    How many of the following statements regarding wrong-way risk and right way risk are correct?
    \begin{enumerate}
        \item Co-movement in risk exposure and default probability producing a decline in overall risk is an example of wrong-way risk
        \item Co-movement in risk exposure and default probability producing an increase in overall counterparty risk is an example of right-way risk
        \item Co-movement in risk exposure and default probability producing neither a decline nor an increase in the overall counterparty risk is an example of wrong-way risk
        \item Co-movement in risk exposure and default probability producing a decline in risk exposure but an increase in counterparty default probability is an example of right-way risk
    \end{enumerate}
\end{question}

\begin{choices}
    \item None
    \item All
    \item Two
    \item Three
\end{choices}

\begin{correctanswer}
    A
\end{correctanswer}

\begin{explanation}
    \textbf{A选项正确。}
    这个选项正确：
    \begin{itemize}
        \item 所有说法都是错误的
        \item 错路风险是指风险敞口和违约概率正相关，导致整体风险增加
        \item 对路风险是指风险敞口和违约概率负相关，导致整体风险减少
    \end{itemize}

    \textbf{B选项错误。}
    这个选项不正确：
    \begin{itemize}
        \item 所有说法都是错误的
        \item 这个选项错误地认为所有说法都正确
        \item 误解了错路风险和对路风险的定义
    \end{itemize}

    \textbf{C选项错误。}
    这个选项不正确：
    \begin{itemize}
        \item 所有说法都是错误的
        \item 这个选项错误地认为有两个说法正确
        \item 误解了错路风险和对路风险的定义
    \end{itemize}

    \textbf{D选项错误。}
    这个选项不正确：
    \begin{itemize}
        \item 所有说法都是错误的
        \item 这个选项错误地认为有三个说法正确
        \item 误解了错路风险和对路风险的定义
    \end{itemize}
\end{explanation}

\checklist{错路风险是敞口与违约概率正相关，对路风险是负相关}

\begin{question}
    Which of the following statements regarding wrong-way risk and right-way risk is correct?
\end{question}

\begin{choices}
    \item A long put option is subject to wrong-way risk if both risk exposure and counterparty default probability decrease
    \item A long call option experiences right-way risk if the interaction between risk exposure and counterparty default probability produces an overall decline in counterparty risk
    \item Declining local currency can decrease the position gain in a foreign currency transaction, while increasing risk exposure of the counterparty
    \item The 2007-2008 credit crisis provides an example of wrong-way risk from the perspective of a long who had sold credit default swaps (CDSs) as protection against bond issuers' default
\end{choices}

\begin{correctanswer}
    B
\end{correctanswer}

\begin{explanation}
    \textbf{A选项错误。}
    这个说法不正确：
    \begin{itemize}
        \item 看跌期权多头在风险敞口和违约概率都增加时才会面临错路风险
        \item 这个选项误解了错路风险的定义
        \item 这个选项描述的是对路风险的情况
    \end{itemize}

    \textbf{B选项正确。}
    这个说法正确：
    \begin{itemize}
        \item 看涨期权多头在风险敞口和违约概率的相互作用导致整体风险下降时
        \item 会经历对路风险
        \item 这符合对路风险的定义
    \end{itemize}

    \textbf{C选项错误。}
    这个说法不正确：
    \begin{itemize}
        \item 本币贬值会增加外汇交易的收益
        \item 同时增加交易对手的风险敞口
        \item 这个选项误解了汇率变化的影响
    \end{itemize}

    \textbf{D选项错误。}
    这个说法不正确：
    \begin{itemize}
        \item 2007-2008年信用危机中的错路风险例子是买入CDS
        \item 而不是卖出CDS
        \item 这个选项误解了历史案例
    \end{itemize}
\end{explanation}

\checklist{看涨期权多头在风险敞口与违约概率负相关时经历对路风险}


\begin{question}
    A risk manager at Bank YENO is describing to a newly hired risk analyst how to estimate and interpret bilateral CVA(BCVA) assuming both Bank YENO, the party making the calculation, and its counterparty can default. To illustrate the estimation of BCVA from Bank YENO's perspective, the manager considers a scenario under which Bank YENO has only one counterparty, Bank GULU, during 2021 and 2022. The manager presents the following results for 2021 and expectations for 2022:
    \begin{itemize}
        \item Bank YENO's CVA estimate in 2021 was USD-13,888
        \item Bank YENO's debt value adjustment(DVA) estimate in 2021 was USD 28,570
        \item The expected positive exposure (EPE) of Bank YENO's derivative position with Bank GULU was USD 35million in 2021, and is expected to remain unchanged in 2022
        \item Bank YENO's expected negative exposure(ENE) with Bank GULU was USD-24 million in 2021, and is expected to remain unchanged in 2022
    \end{itemize}

    The manager also obtains the following information about the two banks:
    \begin{table}[H]
        \centering
        \renewcommand{\arraystretch}{1.2}
        \begin{tabular}{|c|c|c|c|c|}
            \hline
            & \multicolumn{2}{c|}{\textbf{Bank YENO}} & \multicolumn{2}{c|}{\textbf{Bank GULH}} \\
            \hline
            & \textbf{2021} & \textbf{2022} & \textbf{2021} & \textbf{2022} \\
            \hline
            Recovery rate & 85.0\% & 85.0\% & 95.0\% & 95.0\% \\
            \hline
            Probability of default & 1.2\% & 1.2\% & 0.8\% & 0.8\% \\
            \hline
        \end{tabular}
    \end{table}

    In analyzing Bank YENO's BCVA, which of the following statements would the manager be correct to make?
\end{question}

\begin{choices}
    \item The BCVA of Bank YENO in 2022 will equal USD 14,682, implying that the credit quality of Bank YENO has deteriorated compared to 2021
    \item The BCVA of Bank YENO in 2022 will equal USD-53,011, implying that the credit quality of Bank YENO has improved compared to 2021
    \item The BCVA of Bank YENO in 2022 will equal USD-32,141, implying that the credit quality of Bank YENO has improved compared to 2021
    \item The BCVA of Bank YENO in 2022 will equal USD 29,022, implying that the credit quality of Bank YENO has deteriorated compared to 2021
\end{choices}

\begin{correctanswer}
    B
\end{correctanswer}

\begin{explanation}
    \textbf{A选项错误。}
    这个说法不正确：
    \begin{itemize}
        \item BCVA计算错误
        \item 信用质量变化判断错误
        \item 这个选项误解了BCVA的计算和含义
    \end{itemize}

    \textbf{B选项正确。}
    这个说法正确：
    \begin{itemize}
        \item BCVA = CVA - DVA = -13,888 - 28,570 = -53,011
        \item 负的BCVA表示信用质量改善
        \item 这个选项正确计算了BCVA并解释了其含义
    \end{itemize}

    \textbf{C选项错误。}
    这个说法不正确：
    \begin{itemize}
        \item BCVA计算错误
        \item 虽然信用质量变化判断正确
        \item 这个选项误解了BCVA的计算方法
    \end{itemize}

    \textbf{D选项错误。}
    这个说法不正确：
    \begin{itemize}
        \item BCVA计算错误
        \item 信用质量变化判断错误
        \item 这个选项完全误解了BCVA的计算和含义
    \end{itemize}
\end{explanation}

\checklist{BCVA为负表示信用质量改善}




\newpage



\section{考点8:压力测试}
\setcounter{question}{0}
\begin{question}
    Which of the following exposure measures reflects the average distribution of exposures at a specific future date prior to the maturity of the longest maturity transaction within a netting set?
\end{question}

\begin{choices}
    \item Peak exposure
    \item Current exposure
    \item Expected exposure
    \item Expected positive exposure
\end{choices}

\begin{correctanswer}
    C
\end{correctanswer}

\begin{explanation}
    \textbf{A选项错误。}
    这个选项不正确：
    \begin{itemize}
        \item 峰值敞口（Peak exposure）衡量的是极端情况下的敞口
        \item 不是平均分布
        \item 这个选项误解了峰值敞口的定义
    \end{itemize}

    \textbf{B选项错误。}
    这个选项不正确：
    \begin{itemize}
        \item 当前敞口（Current exposure）衡量的是当前的敞口
        \item 不是未来日期的平均分布
        \item 这个选项误解了当前敞口的定义
    \end{itemize}

    \textbf{C选项正确。}
    这个选项正确：
    \begin{itemize}
        \item 期望敞口（Expected exposure）衡量的是特定未来日期的平均敞口分布
        \item 时间范围在净额组合中最长期限之前
        \item 这是期望敞口的定义
    \end{itemize}

    \textbf{D选项错误。}
    这个选项不正确：
    \begin{itemize}
        \item 期望正敞口（Expected positive exposure）是整个期间的平均值
        \item 不是特定未来日期的分布
        \item 这个选项误解了期望正敞口的定义
    \end{itemize}
\end{explanation}

\checklist{期望敞口衡量特定未来日期的平均敞口分布}




\begin{question}
    Which of the following statements best reflects the reason why a financial institution does not need to consider aggregating stresses to the expected positive exposure (EPE) with its loan portfolio?
\end{question}

\begin{choices}
    \item Loans are not sensitive to market variables
    \item Stresses to EPE are not sensitive to market variables
    \item The EPE and the loan portfolio are negatively correlated
    \item The EPE and the loan portfolio are positively correlated
\end{choices}

\begin{correctanswer}
    A
\end{correctanswer}

\begin{explanation}
    \textbf{A选项正确。}
    这个说法正确：
    \begin{itemize}
        \item 贷款对市场变量不敏感
        \item 因此不会因市场变量变化而产生风险敞口变化
        \item 这是不需要考虑EPE压力测试与贷款组合结合的主要原因
    \end{itemize}

    \textbf{B选项错误。}
    这个说法不正确：
    \begin{itemize}
        \item EPE的压力测试对市场变量是敏感的
        \item 这个选项误解了EPE压力测试的特征
        \item 这个选项与实际情况相反
    \end{itemize}

    \textbf{C选项错误。}
    这个说法不正确：
    \begin{itemize}
        \item 相关性不是主要考虑因素
        \item 关键是贷款对市场变量的不敏感性
        \item 这个选项误解了问题的核心
    \end{itemize}

    \textbf{D选项错误。}
    这个说法不正确：
    \begin{itemize}
        \item 相关性不是主要考虑因素
        \item 关键是贷款对市场变量的不敏感性
        \item 这个选项误解了问题的核心
    \end{itemize}
\end{explanation}

\checklist{贷款对市场变量不敏感，无需考虑EPE压力测试}


\begin{question}
    Is the following statement on the treatment of counterparty credit risk (CCR) correct?
    "Treating CCR as a market risk does not allow an institution to hedge market risk losses, and it exposes the institution to declines in counterparty creditworthiness and default."
\end{question}

\begin{choices}
    \item The statement is correct with regard to both hedging market risk losses and counterparty creditworthiness and default
    \item The statement is incorrect with regard to both hedging market risk losses and counterparty creditworthiness and default
    \item The statement is correct with regard to hedging market risk losses only
    \item The statement is correct with regard to counterparty creditworthiness and default only
\end{choices}

\begin{correctanswer}
    D
\end{correctanswer}

\begin{explanation}
    \textbf{A选项错误。}
    这个说法不正确：
    \begin{itemize}
        \item 将CCR视为市场风险允许对冲市场风险损失
        \item 这个选项误解了CCR对冲的可能性
        \item 这个选项完全错误
    \end{itemize}

    \textbf{B选项错误。}
    这个说法不正确：
    \begin{itemize}
        \item 将CCR视为市场风险确实会暴露于交易对手信用质量下降和违约风险
        \item 这个选项误解了CCR的风险特征
        \item 这个选项完全错误
    \end{itemize}

    \textbf{C选项错误。}
    这个说法不正确：
    \begin{itemize}
        \item 将CCR视为市场风险允许对冲市场风险损失
        \item 这个选项误解了CCR对冲的可能性
        \item 这个选项部分错误
    \end{itemize}

    \textbf{D选项正确。}
    这个说法正确：
    \begin{itemize}
        \item 将CCR视为市场风险确实会暴露于交易对手信用质量下降和违约风险
        \item 但允许对冲市场风险损失
        \item 这个选项准确描述了CCR的风险特征
    \end{itemize}
\end{explanation}

\checklist{将CCR视为市场风险可对冲市场风险，但会暴露于对手方违约风险}



\begin{question}
    An analyst notes that stress testing current exposure is problematic because aggregating results is typically not meaningful, although it is easy to account for the credit quality of the counterparty. Are the analyst's statements correct?
\end{question}

\begin{choices}
    \item The analyst is correct with regard to both aggregating results and credit quality
    \item The analyst is correct with regard to aggregating results only
    \item The analyst is correct with regard to credit quality only
    \item The analyst is incorrect with regard to both aggregating results and credit quality
\end{choices}

\begin{correctanswer}
    B
\end{correctanswer}

\begin{explanation}
    \textbf{A选项错误。}
    这个说法不正确：
    \begin{itemize}
        \item 分析师关于汇总结果的说法是正确的
        \item 但关于信用质量的说法是错误的
        \item 这个选项错误地认为两个说法都正确
    \end{itemize}

    \textbf{B选项正确。}
    这个说法正确：
    \begin{itemize}
        \item 汇总压力测试结果确实没有意义
        \item 因为同时违约的可能性很小
        \item 但压力测试并不考虑交易对手的信用质量
    \end{itemize}

    \textbf{C选项错误。}
    这个说法不正确：
    \begin{itemize}
        \item 分析师关于信用质量的说法是错误的
        \item 压力测试不考虑交易对手的信用质量
        \item 这个选项错误地认为信用质量说法正确
    \end{itemize}

    \textbf{D选项错误。}
    这个说法不正确：
    \begin{itemize}
        \item 分析师关于汇总结果的说法是正确的
        \item 这个选项错误地认为两个说法都错误
        \item 这个选项完全误解了分析师的陈述
    \end{itemize}
\end{explanation}

\checklist{压力测试汇总结果无意义，且不考虑对手方信用质量}



\begin{question}
    Which of the following statements regarding exposure measures is not accurate?
\end{question}

\begin{choices}
    \item Assuming no recovery in bankruptcy, current exposure is the greater of zero and the market value of a transaction or portfolio of transactions that the investor would be lost when a counterparty defaulted
    \item Potential Future Exposure measures the distribution of exposures of low percentile
    \item Expected exposure measures the mean distribution of exposures at a given future date prior to the maturity of the longest maturity exposure in the netting group
    \item The weights of expected positive exposure could represent the proportion that an individual expected exposures of the entire time interval
\end{choices}

\begin{correctanswer}
    B
\end{correctanswer}

\begin{explanation}
    \textbf{A选项错误。}
    这个说法是正确的：
    \begin{itemize}
        \item 当前敞口是在交易对手违约时投资者可能损失的市场价值
        \item 假设破产时没有回收
        \item 取零和市场价值的较大值
    \end{itemize}

    \textbf{B选项正确。}
    这个说法是错误的：
    \begin{itemize}
        \item 潜在未来敞口（PFE）衡量的是高百分位数的敞口分布
        \item 而不是低百分位数的敞口分布
        \item 这个选项误解了PFE的定义
    \end{itemize}

    \textbf{C选项错误。}
    这个说法是正确的：
    \begin{itemize}
        \item 期望敞口衡量的是在给定未来日期的敞口均值分布
        \item 时间范围在净额组合中最长期限之前
        \item 这个选项准确描述了期望敞口的定义
    \end{itemize}

    \textbf{D选项错误。}
    这个说法是正确的：
    \begin{itemize}
        \item 期望正敞口的权重可以代表单个期望敞口在整个时间区间中的比例
        \item 这是期望正敞口计算的基本原理
        \item 这个选项准确描述了期望正敞口的权重含义
    \end{itemize}
\end{explanation}

\checklist{PFE衡量高百分位数敞口分布，而非低百分位数}



\begin{question}
    Which of the following statements best reflects the reason why a financial institution does not need to consider aggregating stresses to the expected positive exposure (EPE) with its loan portfolio?
\end{question}

\begin{choices}
    \item Loans are not sensitive to market variables
    \item Stresses to EPE are not sensitive to market variables
    \item The EPE and the loan portfolio are negatively correlated
    \item The EPE and the loan portfolio are positively correlated
\end{choices}

\begin{correctanswer}
    A
\end{correctanswer}

\begin{explanation}
    \textbf{A选项正确。}
    这个说法正确：
    \begin{itemize}
        \item 贷款对市场变量不敏感
        \item 因此不会因市场变量变化而产生风险敞口变化
        \item 这是不需要考虑EPE压力测试与贷款组合结合的主要原因
    \end{itemize}

    \textbf{B选项错误。}
    这个说法不正确：
    \begin{itemize}
        \item EPE的压力测试对市场变量是敏感的
        \item 这个选项误解了EPE压力测试的特征
        \item 这个选项与实际情况相反
    \end{itemize}

    \textbf{C选项错误。}
    这个说法不正确：
    \begin{itemize}
        \item 相关性不是主要考虑因素
        \item 关键是贷款对市场变量的不敏感性
        \item 这个选项误解了问题的核心
    \end{itemize}

    \textbf{D选项错误。}
    这个说法不正确：
    \begin{itemize}
        \item 相关性不是主要考虑因素
        \item 关键是贷款对市场变量的不敏感性
        \item 这个选项误解了问题的核心
    \end{itemize}
\end{explanation}

\checklist{贷款对市场变量不敏感，无需考虑EPE压力测试}


\begin{question}
    An analyst notes that stress testing current exposure is problematic because aggregating results is typically not meaningful, although it is easy to account for the credit quality of the counterparty. Are the analyst's statements correct?
\end{question}

\begin{choices}
    \item The analyst is correct with regard to both aggregating results and credit quality
    \item The analyst is correct with regard to aggregating results only
    \item The analyst is correct with regard to credit quality only
    \item The analyst is incorrect with regard to both aggregating results and credit quality
\end{choices}

\begin{correctanswer}
    B
\end{correctanswer}

\begin{explanation}
    \textbf{A选项错误。}
    这个说法不正确：
    \begin{itemize}
        \item 分析师关于汇总结果的说法是正确的
        \item 但关于信用质量的说法是错误的
        \item 这个选项错误地认为两个说法都正确
    \end{itemize}

    \textbf{B选项正确。}
    这个说法正确：
    \begin{itemize}
        \item 汇总压力测试结果确实没有意义
        \item 因为同时违约的可能性很小
        \item 但压力测试并不考虑交易对手的信用质量
    \end{itemize}

    \textbf{C选项错误。}
    这个说法不正确：
    \begin{itemize}
        \item 分析师关于信用质量的说法是错误的
        \item 压力测试不考虑交易对手的信用质量
        \item 这个选项错误地认为信用质量说法正确
    \end{itemize}

    \textbf{D选项错误。}
    这个说法不正确：
    \begin{itemize}
        \item 分析师关于汇总结果的说法是正确的
        \item 这个选项错误地认为两个说法都错误
        \item 这个选项完全误解了分析师的陈述
    \end{itemize}
\end{explanation}

\checklist{压力测试汇总结果无意义，且不考虑对手方信用质量}


\begin{question}
    An analyst notes that stress testing current exposure is problematic because aggregating results is typically not meaningful, although it is easy to account for the credit quality of the counterparty. Are the analyst's statements correct?
\end{question}

\begin{choices}
    \item The analyst is correct with regard to both aggregating results and credit quality
    \item The analyst is correct with regard to aggregating results only
    \item The analyst is correct with regard to credit quality only
    \item The analyst is incorrect with regard to both aggregating results and credit quality
\end{choices}

\begin{correctanswer}
    B
\end{correctanswer}

\begin{explanation}
    \textbf{A选项错误。}
    这个说法不正确：
    \begin{itemize}
        \item 分析师关于汇总结果的说法是正确的
        \item 但关于信用质量的说法是错误的
        \item 这个选项错误地认为两个说法都正确
    \end{itemize}

    \textbf{B选项正确。}
    这个说法正确：
    \begin{itemize}
        \item 汇总压力测试结果确实没有意义
        \item 因为同时违约的可能性很小
        \item 但压力测试并不考虑交易对手的信用质量
    \end{itemize}

    \textbf{C选项错误。}
    这个说法不正确：
    \begin{itemize}
        \item 分析师关于信用质量的说法是错误的
        \item 压力测试不考虑交易对手的信用质量
        \item 这个选项错误地认为信用质量说法正确
    \end{itemize}

    \textbf{D选项错误。}
    这个说法不正确：
    \begin{itemize}
        \item 分析师关于汇总结果的说法是正确的
        \item 这个选项错误地认为两个说法都错误
        \item 这个选项完全误解了分析师的陈述
    \end{itemize}
\end{explanation}

\checklist{压力测试汇总结果无意义，且不考虑对手方信用质量}

\begin{question}
    Is the following statement on the treatment of counterparty credit risk (CCR) correct?
    "Treating CCR as a market risk does not allow an institution to hedge market risk losses, and it exposes the institution to declines in counterparty creditworthiness and default."
\end{question}

\begin{choices}
    \item The statement is correct with regard to both hedging market risk losses and counterparty creditworthiness and default
    \item The statement is incorrect with regard to both hedging market risk losses and counterparty creditworthiness and default
    \item The statement is correct with regard to hedging market risk losses only
    \item The statement is correct with regard to counterparty creditworthiness and default only
\end{choices}

\begin{correctanswer}
    D
\end{correctanswer}

\begin{explanation}
    \textbf{A选项错误。}
    这个说法不正确：
    \begin{itemize}
        \item 将CCR视为市场风险允许对冲市场风险损失
        \item 这个选项误解了CCR对冲的可能性
        \item 这个选项完全错误
    \end{itemize}

    \textbf{B选项错误。}
    这个说法不正确：
    \begin{itemize}
        \item 将CCR视为市场风险确实会暴露于交易对手信用质量下降和违约风险
        \item 这个选项误解了CCR的风险特征
        \item 这个选项完全错误
    \end{itemize}

    \textbf{C选项错误。}
    这个说法不正确：
    \begin{itemize}
        \item 将CCR视为市场风险允许对冲市场风险损失
        \item 这个选项误解了CCR对冲的可能性
        \item 这个选项部分错误
    \end{itemize}

    \textbf{D选项正确。}
    这个说法正确：
    \begin{itemize}
        \item 将CCR视为市场风险确实会暴露于交易对手信用质量下降和违约风险
        \item 但允许对冲市场风险损失
        \item 这个选项准确描述了CCR的风险特征
    \end{itemize}
\end{explanation}

\checklist{将CCR视为市场风险可对冲市场风险，但会暴露于对手方违约风险}


\begin{question}
    Risk managers at a commercial bank are preparing a stress test of the bank's capital ratios that assumes a severe financial market shock over a 2-year time horizon. The managers want to ensure that the stress test appropriately captures the impact of the stress scenario on the bank's investment portfolios, as well as on the bank's balance sheet items over the stress test horizon. Which of the following would be the most appropriate assumption or methodology for the managers to include in the stress test?
\end{question}

\begin{choices}
    \item Yields on speculative-grade corporate debt and short-term AAA-rated sovereign debt increase simultaneously
    \item The default probabilities of derivatives counterparties increase as the bank's exposures to those counterparties increase
    \item The volume of the bank's used auto loans falls much more than the volume of its new auto loans
    \item The stress test is calibrated so that the bank is assumed to operate at the minimum capital ratio required to remain solvent
\end{choices}

\begin{correctanswer}
    B
\end{correctanswer}

\begin{explanation}
    \textbf{A选项错误。}
    这个说法不正确：
    \begin{itemize}
        \item 投机级公司债和AAA级主权债收益率不会同时上升
        \item 在压力情景下，投资者会转向高质量资产
        \item 这个选项不符合市场行为特征
    \end{itemize}

    \textbf{B选项正确。}
    这个说法正确：
    \begin{itemize}
        \item 衍生品交易对手的违约概率会随着敞口增加而上升
        \item 这反映了压力情景下的信用风险传导
        \item 这个选项准确描述了压力测试中的风险特征
    \end{itemize}

    \textbf{C选项错误。}
    这个说法不正确：
    \begin{itemize}
        \item 二手车贷款和新车贷款的下降幅度差异不是关键因素
        \item 这个选项过于关注具体业务细节
        \item 这个选项不符合压力测试的整体性要求
    \end{itemize}

    \textbf{D选项错误。}
    这个说法不正确：
    \begin{itemize}
        \item 压力测试不应该假设银行仅维持最低资本要求
        \item 这会导致测试结果过于乐观
        \item 这个选项不符合审慎监管要求
    \end{itemize}
\end{explanation}

\checklist{压力测试中对手方违约概率随敞口增加而上升}


\begin{question}

\end{question}

\begin{choices}
    \item 
\end{choices}

\begin{correctanswer}
    B
\end{correctanswer}

\begin{explanation}
    \textbf{A选项错误。}
    \begin{itemize}
        \item 
    \end{itemize}
    \textbf{B选项错误。}
    \begin{itemize}
        \item 
    \end{itemize}
    \textbf{C选项错误。}
    \begin{itemize}
        \item 
    \end{itemize}
    \textbf{D选项错误。}
    \begin{itemize}
        \item 
    \end{itemize}
\end{explanation}
\checklist{}



\begin{question}

\end{question}

\begin{choices}
    \item 
\end{choices}

\begin{correctanswer}
    B
\end{correctanswer}

\begin{explanation}
    \textbf{A选项错误。}
    \begin{itemize}
        \item 
    \end{itemize}
    \textbf{B选项错误。}
    \begin{itemize}
        \item 
    \end{itemize}
    \textbf{C选项错误。}
    \begin{itemize}
        \item 
    \end{itemize}
    \textbf{D选项错误。}
    \begin{itemize}
        \item 
    \end{itemize}
\end{explanation}
\checklist{}


\newpage





\end{document}